% संपूर्ण मराठी ग्रंथातील प्रकरण 29: $\Th{Q}$ मधील निरूपणीयता
% हा संपादनयोग्य स्रोतखंड आहे. संकलनासाठी संपूर्ण स्रोतसंग्रहातील
% अक्षररूपे, पूर्वभाग, चिन्हव्याख्या आणि आकृती-स्रोत आवश्यक आहेत.
% मूळ: Open Logic Project; मजकूर आणि रूपांतर CC BY 4.0.
% यंत्रानुवाद, दुरुस्ती आणि तपासणी: OpenAI Codex —
% GPT-5.6 Sol आणि GPT-6 Sol, दोन्ही Ultra effort.
% दुरुस्ती-पुनरावलोकन: GPT-6.1 Sol, Ultra effort.
\chapter{$\Th{Q}$ मधील निरूपणीयता}\label{inc:req::chap}








\section{परिचय}\label{inc:req:int:sec}

अपूर्णता प्रमेये अशा उपपत्तींना लागू होतात ज्यांत संगणनक्षम
फलनांविषयीची मूलभूत तथ्ये व्यक्त आणि सिद्ध करता येतात. अशा प्रकारची
अगदी किमान उपपत्ती आपण पाहणार आहोत; तिला ``$\Th{Q}$'' म्हणतात
(राफेल रॉबिन्सन यांच्या नावावरून तिला कधीकधी ``रॉबिन्सनची $Q$''
असेही म्हणतात). एखादे फलन $\Th{Q}$ मध्ये \emph{निरूपणीय} असणे म्हणजे
काय, हे आपण सांगू आणि मग पुढील विधान सिद्ध करू:
\begin{quote}
  फलन $\Th{Q}$ मध्ये निरूपणीय असते, तर आणि तरच ते संगणनक्षम असते.
\end{quote}
यामुळे आपल्याला संगणनक्षमतेचे आणखी एक प्रतिमान मिळते. शिवाय,
थांबण्याच्या समस्येचे त्याकडे न्यूनीकरण करून
$\Setabs{\varphi}{\Th{Q} \Proves \varphi}$ हा संच निर्णेय नाही, हे दाखवण्यासाठीही
त्याचा उपयोग करू. पुढे याहून अधिक बलवान निष्कर्ष आपण सिद्ध करू.

$\Th{Q}$ ची भाषा ही अंकगणिताची भाषा आहे; $\Th{Q}$ मध्ये पुढील
स्वयंसिद्धके आहेत. ती एकरूपता असलेल्या प्रथम-क्रम तर्कशास्त्रातील
इतर स्वयंसिद्धके आणि नियम यांच्यासोबत वापरायची आहेत:
\begin{align*}
& \lforall[x][\lforall[y][(\eq[x'][y'] \lif \eq[x][y])]] \tag{$\mathsf{Q}_1$}\\
& \lforall[x][\eq/[\Obj 0][x']] \tag{$\mathsf{Q}_2$}\\
& \lforall[x][(\eq[x][\Obj 0] \lor \lexists[y][\eq[x][y']])] \tag{$\mathsf{Q}_3$}\\
& \lforall[x][\eq[(x + \Obj 0)][x]] \tag{$\mathsf{Q}_4$}\\
& \lforall[x][\lforall[y][\eq[(x + y')][(x + y)']]] \tag{$\mathsf{Q}_5$}\\
& \lforall[x][\eq[(x \times \Obj 0)][\Obj 0]] \tag{$\mathsf{Q}_6$}\\
& \lforall[x][\lforall[y][\eq[(x \times y')][((x \times y) + x)]]] \tag{$\mathsf{Q}_7$}\\
& \lforall[x][\lforall[y][(x < y \liff \lexists[z][\eq[(z' + x)][y]])]] \tag{$\mathsf{Q}_8$}
\end{align*}
प्रत्येक नैसर्गिक संख्या $n$ साठी $\num{n}$ हा संख्यांक
$\Obj{0}^{\prime\prime\ldots\prime}$ हे पद म्हणून परिभाषित करा;
त्यात एकूण $n$ उत्तरवर्ती चिन्हे असतात. म्हणजे $\num{0}$ हा
स्थिरांक~$\Obj{0}$ एकटाच, $\num{1}$ हे $\Obj{0}'$,
$\num{2}$ हे $\Obj{0}''$, इत्यादी.

अंकगणिताची उपपत्ती म्हणून $\Th{Q}$ \emph{अत्यंत} दुर्बल आहे;
उदाहरणार्थ, $\lforall[x][\eq/[x][x']]$ किंवा
$\lforall[x][\lforall[y][(x + y) = (y + x)]]$ यांसारखी
अगदी साधी तथ्येसुद्धा त्यात सिद्ध करता येत नाहीत. पण $\Th{Q}$
इतकी रंजक असण्याचे मोठे कारण ती \emph{दुर्बल} आहे, हेच आहे.
अपूर्णता प्रमेय लागू होण्याइतकीच तिची शक्ती आहे. $\Th{Q}$
रंजक असण्याचे आणखी एक कारण म्हणजे तिचा स्वयंसिद्धकांचा संच
\emph{सांत} आहे.

$\Th{Q}$ पेक्षा अधिक बलवान उपपत्ती \emph{पेआनो अंकगणित}
$\Th{PA}$ हिला $\Th{Q}$ मध्ये विगमन-रूपबंध जोडून मिळवतात:
\[(\varphi(\Obj 0) \land \lforall[x][(\varphi(x) \lif \varphi(x'))]) \lif \lforall[x][\varphi(x)]
\]
येथे $\varphi(x)$ हे कोणतेही सूत्र असते. $\varphi(x)$ मध्ये $x$ व्यतिरिक्त
इतर चले मुक्त असतील, तर त्या सर्वांना बांधण्यासाठी
आरंभी सार्वत्रिक संख्यक जोडतो (म्हणजे विगमन-रूपबंधाचे संबंधित
उदाहरण वाक्य असेल). उदाहरणार्थ, $\varphi(x, y)$ मध्ये
चल~$y$ हेसुद्धा मुक्त असेल, तर संबंधित उदाहरण असे आहे:
\[\lforall[y][((\varphi(\Obj 0) \land \lforall[x][(\varphi(x) \lif \varphi(x'))]) \lif
  \lforall[x][\varphi(x)])]
\]
विगमन-रूपबंधाची उदाहरणे वापरून $\Th{Q}$ च्या स्वयंसिद्धकांपेक्षा
$\Th{PA}$ च्या स्वयंसिद्धकांपासून बरेच अधिक सिद्ध करता येते.
खरे तर, पेआनो अंकगणितात सिद्ध करता येत नाहीत अशी नैसर्गिक
संख्यांविषयीची ``स्वाभाविक'' विधाने शोधणे बरेच कठीण असते!{}

\begin{defn}
\label{inc:req:int:defn:representable-fn}
  नैसर्गिक संख्यांपासून नैसर्गिक संख्यांकडे जाणारे
  $f(x_0,\ldots,x_k)$ हे फलन $\Th{Q}$ मध्ये
  {\em निरूपणीय} आहे असे म्हणतात, जर असे सूत्र
  $\varphi_f(x_0,\dots,x_k,y)$ असेल की
  $f(n_0,\dots,n_k) = m$ असेल तेव्हा $\Th{Q}$ पुढील दोन्ही
  विधाने सिद्ध करते:
\begin{enumerate}
\item\label{inc:req:int:defn:rep:a} $\varphi_f(\num{n_0}, \dots, \num{n_k}, \num{m})$
\item\label{inc:req:int:defn:rep:b} $\lforall[y][(\varphi_f(\num{n_0}, \dots,
\num{n_k}, y) \lif \num{m} = y)]$.
\end{enumerate}
\end{defn}

ही व्याख्या इतर मार्गांनीही मांडता येते. उदाहरणार्थ, $\Th{Q}$
$\lforall[y][(\varphi_f(\num{n_0}, \dots,
\num{n_k}, y) \liff \eq[y][\num{m}])]$ सिद्ध करते, अशी
समतुल्य अटही आपण ठेवू शकलो असतो.

\begin{thm}
\label{inc:req:int:thm:representable-iff-comp}
फलन $\Th{Q}$ मध्ये निरूपणीय असते, तर आणि तरच ते संगणनक्षम असते.
\end{thm}

या प्रमेयाच्या पुराव्याला दोन दिशा आहेत. विन्यासमीमांसेचे अंकगणितीकरण
झाल्यावर डावीकडून उजवीकडे जाणारी दिशा तुलनेने सोपी आहे. दुसऱ्या
दिशेसाठी अधिक काम लागते. मूलभूत कल्पना अशी: ``संगणनक्षम'' याचा
अर्थ नेमका करण्यासाठी ``सामान्य पुनरावर्ती'' हा पर्याय निवडतो आणि
प्रत्येक सामान्य पुनरावर्ती फलन $\Th{Q}$ मध्ये निरूपणीय आहे,
हे दाखवतो. लक्षात घ्या, शून्य फलन~$\Zero$, उत्तरवर्ती फलन~$\Succ$
आणि प्रक्षेपण फलने~$\Proj{n}{i}$ यांच्यापासून संयोजन, आदिम पुनरावर्तन
आणि नियमित लघुतमीकरण वापरून परिभाषित करता येणारे फलन म्हणजे
सामान्य पुनरावर्ती फलन. म्हणून प्रत्येक सामान्य पुनरावर्ती फलन
$\Th{Q}$ मध्ये निरूपणीय आहे हे दाखवण्याचा एक मार्ग असा: मूलभूत
फलने निरूपणीय आहेत हे दाखवायचे आणि काही फलने निरूपणीय असताना
त्यांच्यापासून संयोजन, आदिम पुनरावर्तन किंवा नियमित लघुतमीकरणाने
बनणारी फलनेही निरूपणीय असतात, हे दाखवायचे. दुसऱ्या शब्दांत,
मूलभूत फलने निरूपणीय आहेत आणि निरूपणीय फलनांचा वर्ग संयोजन,
आदिम पुनरावर्तन व नियमित लघुतमीकरण यांखाली ``बंद'' आहे,
हे दाखवता येईल. त्यावरून प्रत्येक सामान्य पुनरावर्ती फलन
निरूपणीय असल्याची हमी मिळते.

पण निरूपणीय फलनांचा वर्ग आदिम पुनरावर्तनाखाली बंद आहे हे दाखवण्याची
पायरी कठीण ठरते. ती टाळण्यासाठी, आधी आपण दाखवतो की प्रत्यक्षात
आदिम पुनरावर्तनाशिवायही काम चालते. म्हणजे, प्रत्येक सामान्य
पुनरावर्ती फलन मूलभूत फलनांपासून केवळ संयोजन आणि नियमित
लघुतमीकरण वापरून परिभाषित करता येते, हे दाखवतो. त्यासाठी आदिम
पुनरावर्तन एका विशिष्ट नियमित लघुतमीकरणाने करता येते, हे दाखवतो.
मात्र हे साधण्यासाठी काही आणखी मूलभूत फलने घ्यावी लागतात:
बेरीज, गुणाकार आणि एकरूपता संबंधाचे जातिबोधक फल~$\Char{=}$. मग
या \emph{सर्व} मूलभूत फलनांची $\Th{Q}$ मधील निरूपणीयता आणि
निरूपणीय फलनांचा संयोजन व नियमित लघुतमीकरण यांखालील बंदपणा
दाखवून हे प्रमेय सिद्ध करता येते.











\section{$\Th{Q}$ मध्ये निरूपणीय फलने संगणनक्षम असतात}\label{inc:req:rpc:sec}

$\Th{Q}$ मध्ये निरूपणीय असलेले प्रत्येक फलन संगणनक्षम आहे,
हे आपण सिद्ध करू. त्याआधी $\Th{Q}$ मध्ये निरूपणीय फलनांविषयी
एक साहाय्यक विधान सिद्ध करावे लागेल.

\begin{lem}\label{inc:req:rpc:lem:rep-q}
  $f(x_0,
  \dots, x_k)$ हे $\Th{Q}$ मध्ये निरूपणीय असेल, तर
  असे सूत्र~$\varphi_f(x_0, \dots, x_k, y)$ असते की
  \[  \Th{Q} \Proves \varphi_f(\num{n_0}, \dots, \num{n_k}, \num{m})
  \quad\text{तेव्हा आणि केवळ तेव्हाच}\quad m = f(n_0, \dots, n_k).
  \]
\end{lem}

\begin{proof}
  ``तर'' हा भाग
\ref{inc:req:int:defn:representable-fn}\ref{inc:req:int:defn:rep:a} मधून मिळतो.
``तरच'' हा भाग पुढीलप्रमाणे दिसतो: समजा $\Th{Q} \Proves \varphi_f(\num{n_0},
\dots, \num{n_k}, \num{m})$, पण $m \neq f(n_0, \dots, n_k)$.
$l = f(n_0, \dots, n_k)$ असे धरू.
\ref{inc:req:int:defn:representable-fn}\ref{inc:req:int:defn:rep:a} नुसार $\Th{Q}
\Proves \varphi_f(\num{n_0}, \dots, \num{n_k}, \num{l})$.
\ref{inc:req:int:defn:representable-fn}\ref{inc:req:int:defn:rep:b} नुसार
$\Th{Q} \Proves \lforall[y][(\varphi_f(\num{n_0}, \dots, \num{n_k}, y) \lif \num{l} =
y)]$. तर्कशास्त्र वापरून आणि $\Th{Q} \Proves
\varphi_f(\num{n_0}, \dots, \num{n_k}, \num{m})$ या गृहीतकावरून
$\Th{Q} \Proves \eq[\num{l}][\num{m}]$ मिळते. दुसरीकडे,
\ref{inc:req:bre:lem:q-proves-neq} नुसार $\Th{Q} \Proves
\eq/[\num{l}][\num{m}]$. म्हणून $\Th{Q}$ विसंगत ठरेल.
पण ते अशक्य आहे: मानक प्रतिमान $\Th{Q}$ ची पूर्ती करते
(पाहा \ref{inc:int:def:def:standard-model}), म्हणजे
$\Sat{N}{\Th{Q}}$; आणि निर्दोषता प्रमेयानुसार प्रत्येक
पूर्ततायोग्य उपपत्ती सुसंगत असते
(\ref{fol:axd:sou:cor:consistency-soundness}, \ref{fol:seq:sou:cor:consistency-soundness}, \ref{fol:ntd:sou:cor:consistency-soundness}, \ref{fol:tab:sou:cor:consistency-soundness}).
\end{proof}

\begin{lem}
$\Th{Q}$ मध्ये निरूपणीय असलेले प्रत्येक फलन संगणनक्षम आहे.
\end{lem}

\begin{proof}
हे खरे असण्यामागील अंतर्ज्ञानी कल्पना आधी पाहू. $f$ चे मूल्य
काढण्यासाठी पुढील पद्धत वापरू. अंकगणिताच्या भाषेतील सर्व संभाव्य
निष्पत्त्या~$\delta$ यांची यादी करायची; हे यांत्रिकरीत्या
करता येते. प्रत्येकासाठी, ती $\varphi_f(\num{n_0}, \dots,
\num{n_k}, \num{m})$ या रूपाच्या सूत्राची
निष्पत्ती आहे का ते तपासायचे (हे $f$ चे $\Th{Q}$
मधील निरूपण करणारे सूत्र आहे; पाहा \ref{inc:req:rpc:lem:rep-q}).
असेल, तर \ref{inc:req:rpc:lem:rep-q} नुसार $m = f(n_0, \dots, n_k)$;
म्हणजे $f$ चे मूल्य सापडले. हा शोध संपतो, कारण
$\Th{Q} \Proves \varphi_f(\num{n_0}, \dots, \num{n_k},
\num{f(n_0, \dots, n_k)})$; त्यामुळे योग्य प्रकारची
$\delta$ अखेरीस सापडते.

ही मांडणी अजून पूर्णपणे नेमकी नाही: आपली पद्धत फक्त संख्यांवर
न चालता निष्पत्त्या आणि सूत्रे वर चालते, आणि
``सर्व संभाव्य निष्पत्त्यांची यादी'' यांत्रिकरीत्या कशी
करता येते हे आपण अजून स्पष्ट केलेले नाही. पण आपण पाहिलेच आहे
की पदे, सूत्रे आणि निष्पत्त्या यांना गोडेल संख्यांनी
संकेतबद्ध करता येते. संगणनाचे नेमके प्रतिमान म्हणून सामान्य
पुनरावर्ती फलनेही आपण मांडली आहेत. तसेच
$\Prf[\Th{Q}](d,y)$ हा संबंध (आदिम) पुनरावर्ती आहे:
नैसर्गिक निगमनातील थेट निष्पत्ती-संकेतांकाचा अर्थ घेतल्यास,
$d$ ही $\Th{Q}$ च्या स्वयंसिद्धकांपासून $y$ ही गोडेल संख्या
असलेल्या सूत्राच्या निष्पत्तीची गोडेल संख्या
असेल, तर आणि तरच हा संबंध खरा असतो. आपल्याला आणखी लागणारी आदिम
पुनरावर्ती फलने म्हणजे $\fn{num}$
(\ref{inc:art:trm:prop:num-primrec}) आणि $\fn{Subst}$
(\ref{inc:art:sub:prop:subst-primrec}). यांपासून लघुतमीकरणाने
$f$ परिभाषित करता येते; म्हणून $f$ पुनरावर्ती आहे.

प्रथम पुढील व्याख्या करा:
\begin{multline*}
  A(n_0, \dots, n_k, m) = \\
  \fn{Subst}(\fn{Subst}(\dots\fn{Subst}(\Gn{\varphi_f}, \fn{num}(n_0), \Gn{x_0}),\\
  \dots),
  \fn{num}(n_k),  \Gn{x_k}), \fn{num}(m), \Gn{y})
\end{multline*}
हे गुंतागुंतीचे दिसते, पण हे फक्त $A(n_0, \dots, n_k,
m) = \Gn{\varphi_f(\num{n_0}, \dots, \num{n_k}, \num{m})}$
हे फलन आहे.

आता $R(n_0, \dots, n_k, s)$ हा संबंध पाहा. $(s)_0$ ही
$\varphi_f(\num{n_0}, \dots, \num{n_k}, \num{(s)_1})$ च्या
$\Th{Q}$ पासूनच्या निष्पत्तीची गोडेल संख्या असेल
तेव्हा हा संबंध खरा असतो:
\[R(n_0, \dots, n_k, s) \quad\text{तेव्हा आणि केवळ तेव्हाच}\quad \Prf[\Th{Q}]((s)_0, A(n_0,
\dots, n_k, (s)_1))
\]
$R(n_0, \dots, n_k, s)$ खरा असणारा एखादा $s$ मिळाला,
तर $(s)_0$ आणि $(s)_1$ ही संख्यांची जोडी मिळाली आहे:
$(s)_0$ ही $\varphi_f(\num{n_0}, \dots,
\num{n_k}, \num{(s)_1})$ च्या निष्पत्तीची गोडेल संख्या
असते. म्हणून $s$ शोधणे म्हणजे अनौपचारिक पुराव्यातील $d$ आणि
$m$ ही जोडी शोधण्यासारखे आहे. असा $s$ ``शोधणारे'' संगणनक्षम
फलन नियमित लघुतमीकरणाने परिभाषित करता येते. लक्षात घ्या,
$R$ साठी हा शोध नियमित आहे: प्रत्येक $n_0$, \dots, $n_k$
साठी $\Th{Q} \Proves \varphi_f(\num{n_0}, \dots,
\num{n_k}, \num{f(n_0, \dots, n_k)})$ हे सिद्ध करणारी निष्पत्ती~$\delta$
असते. म्हणून $s = \tuple{\Gn{\delta}, f(n_0, \dots, n_k)}$
घेतल्यास $R(n_0, \dots, n_k, s)$ खरा असतो. त्यामुळे $f$
पुढीलप्रमाणे लिहिता येते:
\[f(n_0,\dots,n_{k}) = (\umin{s}{R(n_0, \dots, n_k, s)})_1.
\]
\end{proof}











\section{बीटा फलनावरील साहाय्यक विधान}\label{inc:req:bet:sec}


बेरीज, गुणाकार आणि $\Char{=}$ उपलब्ध असतील तर आदिम पुनरावर्तन
करता येते हे दाखवण्यासाठी क्रमिका हाताळणारी फलने विकसित करावी लागतील.
(घातांक फलनही उपलब्ध असते, तर हे काम सोपे झाले असते.) आदिम
पुनरावर्तन उपलब्ध असताना ``$n$ वी अविभाज्य संख्या'' यांसारख्या
गोष्टी परिभाषित करून तुलनेने सरळ सांकेतीकरण निवडता येत होते.
पण येथे आदिम पुनरावर्तन उपलब्ध नाही---उलट, लघुतमीकरण वापरून
आदिम पुनरावर्तन करता येते हेच दाखवायचे आहे---म्हणून अधिक
युक्तीने काम करावे लागेल.

\begin{lem}
\label{inc:req:bet:lem:beta}
असे फलन $\beta(d,i)$ आहे की प्रत्येक क्रमिका $a_0$,
\dots,~$a_n$ साठी अशी संख्या~$d$ असते की प्रत्येक $i \le n$
साठी $\beta(d,i) = a_i$. शिवाय, $\beta$ मूलभूत फलनांपासून
केवळ संयोजन आणि नियमित लघुतमीकरण वापरून परिभाषित करता येते.
\end{lem}

$d$ हा क्रमिका $\tuple{a_0, \dots, a_n}$ चा संकेतांक आहे
आणि $\beta(d,i)$ तिचा $i$ वा घटक देतो, असे समजा.
(लक्षात घ्या, या ``सांकेतीकरणात'' आपल्याला परिचित असलेले
अविभाज्य संख्यांच्या घातांवर आधारित सांकेतीकरण \emph{वापरलेले नाही}!)
हे साहाय्यक विधान किमान आवश्यक तेवढेच सांगते: क्रमिका जोडता
येतात, त्यांना घटक शेवटी जोडता येतो किंवा संयोजन व नियमित
लघुतमीकरणाने परिभाषित फलने वापरून $a_0$, \dots,~$a_n$
यांपासून $d$ \emph{काढता} येतो, असेही ते म्हणत नाही. ते फक्त
एवढेच सांगते की प्रत्येक क्रमिकेसाठी काही संकेतांक असतो आणि
त्यातून घटक मिळवणारे एक ``विसांकेतीकरण'' फलन आहे.

$\beta$ हे चिन्हांकन गोडेल यांचे आहे. पुन्हा सांगायचे तर,
या साहाय्यक विधानाच्या पुराव्यातील कठीण भाग म्हणजे वरवर
मर्यादित वाटणारी साधने---केवळ संयोजन आणि लघुतमीकरण---वापरून
योग्य $\beta$ परिभाषित करणे; अर्थात बेरीज, गुणाकार आणि
$\Char{=}$ वापरता येतात. या विधानाचे अनेक पुरावे आहेत; पण
त्यांपैकी एक सुबोध मार्ग आजही गोडेल यांची मूळ पद्धत आहे.
त्या पद्धतीत सुनझीचे प्रमेय (परंपरेने ``चिनी शेष प्रमेय'')
हे संख्या-सैद्धान्तिक तथ्य वापरले जाते.

\begin{defn}
दोन नैसर्गिक संख्या $a$ आणि $b$ या \emph{सापेक्षतः अविभाज्य}
(सहमूळ) असतात, तर आणि तरच त्यांचा महत्तम साधारण
विभाजक~$1$ असतो; म्हणजे त्यांचा दुसरा कोणताही सामाईक विभाजक नसतो.
\end{defn}

\begin{defn}
धन भाजक $c>0$ साठी नैसर्गिक संख्या $a$ आणि $b$ या \emph{भाजक~$c$ समशेष},
$a \equiv b \mod c$, असतात, तर आणि तरच
$c \mid (a-b)$; म्हणजे $a$ आणि $b$ यांना $c$ ने भागल्यावर
शेष समान मिळतो.

\end{defn}

सुनझीचे प्रमेय असे आहे:
\begin{thm}
समजा $x_0$, \dots,~$x_n$ या धन नैसर्गिक संख्या परस्पर जोड्यांनी सापेक्षतः
अविभाज्य आहेत. $y_0$, \dots,~$y_n$ या कोणत्याही संख्या असू द्या.
तर अशी संख्या $z$ आहे की
\begin{align*}
z & \equiv y_0 \mod x_0 \\
z & \equiv y_1 \mod x_1 \\
& \vdots  \\
z & \equiv y_n \mod x_n.
\end{align*}
\end{thm}

सुनझीचे प्रमेय आपण पुढीलप्रमाणे वापरू: $x_0$, \dots,~$x_n$
या अनुक्रमे $y_0$, \dots,~$y_n$ पेक्षा मोठ्या असतील, तर
$z$ हा क्रमिका $\tuple{y_0, \dots,y_n}$ चा संकेतांक घेता येतो.
$y_i$ परत मिळवण्यासाठी $z$ ला $x_i$ ने भागून शेष घ्यायचा.
या सांकेतीकरणासाठी $x_0$, \dots,~$x_n$ यांची योग्य मूल्ये
शोधावी लागतील.

यासाठी पुढील दोन निरीक्षणे उपयोगी पडतील. $y_0$, \dots,~$y_n$
दिली असता, पुढीलप्रमाणे घ्या:
\begin{align*}
j &= \max(n, y_0 + 1, \dots, y_n + 1), \\
m &= \lcm(1,\dots,j),
\end{align*}
दुसऱ्या ओळीत लघुतम साधारण विभाज्य घेतला आहे.
आणि पुढीलप्रमाणे घ्या:
\begin{align*}
x_0 & = 1 + m \\
x_1 & = 1 + 2 \cdot m \\
x_2 & = 1 + 3 \cdot m \\
& \vdots  \\
x_n & = 1 + (n+1) \cdot m
\end{align*}
मग पुढील दोन गोष्टी खऱ्या आहेत:
\begin{enumerate}
\item\label{inc:req:bet:rel-prime} $x_0,\dots,x_n$ या परस्पर जोड्यांनी
  सापेक्षतः अविभाज्य आहेत.
\item\label{inc:req:bet:less} प्रत्येक $i$ साठी $y_i < x_i$.
\end{enumerate}
\ref{inc:req:bet:rel-prime} हे खरे का, ते पाहू. दोन भिन्न निर्देशांक
निवडा. $p$ ही अविभाज्य संख्या
असेल आणि $p \mid x_i$ व $p \mid x_k$ असतील, तर
$p \mid 1 + (i+1) m$ आणि $p \mid 1 + (k+1) m$.
मग $p$ त्यांच्या फरकालाही विभाजित करते:
\[(1 + (i+1)m) - (1+ (k+1)m) = (i-k) m.
\]
$p$ ही $1 + (i+1)m$ ला विभाजित करते, म्हणून ती $m$ ला
विभाजित करू शकत नाही (अन्यथा पहिल्या भागाकारात शेष~$1$
उरेल). त्यामुळे $p$ ही $(i-k)m$ ला विभाजित करत असल्याने
$p$ ही $i-k$ लाही विभाजित करते. पण $\left|i-k\right|$ जास्तीत
जास्त~$n$ आहे, आणि $j \geq n$ असे घेतले आहे; म्हणून
पुन्हा $p \mid m$, हा विरोधाभास. म्हणून $x_i$ आणि
$x_k$ या दोन्हींना विभाजित करणारी अविभाज्य संख्या नाही.
\ref{inc:req:bet:less} ही अट सोपी आहे: $y_i < j \leq m < x_i$.

आता $\beta$ फलनावरील साहाय्यक विधान सिद्ध करू. लक्षात ठेवा,
$0$, उत्तरवर्ती फलन, बेरीज, गुणाकार, $\Char{=}$, प्रक्षेपणे,
आणि नियमित फलनांना लघुतमीकरण लागू करून व संयोजनाने
यांपासून परिभाषित केलेली कोणतीही फलने वापरता येतात.
एखाद्या संबंधाचे जातिबोधक फल असे परिभाषित करता येत असेल,
तर तो संबंधही वापरता येतो. आधीप्रमाणे, हे संबंध बूलीय
संयोग आणि परिबद्ध संख्यकीकरण यांखाली बंद आहेत हे दाखवता येते;
उदाहरणार्थ:
\begin{align*}
\fn{not}(x) & \defis \Char{=}(x,0)\\
\bmin{x \leq z}{R(x,y)} & \defis \umin{x}{(R(x,y) \lor x = z)}\\
\bexists{x \leq z}{R(x,y)} & \defiff R(\bmin{x \leq z}{R(x,y)}, y)
\end{align*}
मग पुढील सर्व गोष्टी आदिम पुनरावर्तनाशिवायही परिभाषित
करता येतात, हे दाखवता येते:
\begin{enumerate}
\item जोडीकरण फलन, $J(x,y) = \frac{1}{2}[(x+y)(x+y+1)] + x$;

\item प्रक्षेपण फलने
\begin{align*}
K(z) & = \bmin{x \leq z}{\bexists{y \leq z}{z = J(x,y)}},\\
L(z) & = \bmin{y \leq z}{\bexists{x \leq z}{z = J(x,y)}};
\end{align*}
\item $x < y$ हा लघुत्वाचा संबंध;
\item $x \mid y$ हा विभाज्यतेचा संबंध;






\item $y$ ला $x$ ने भागल्यावर मिळणारा शेष देणारे
  $\fn{rem}(x,y)$ हे फलन.
\end{enumerate}
शून्य भाजकासाठी या शेष फलाचे मूल्य येथे दिलेले नाही; आधीच्या
शेष फलाच्या परिपाठाप्रमाणे ते शून्य घेता येते. पुढील
वापरातील भाजक मात्र नेहमी धन असतो.
आता पुढील व्याख्या करा:
\begin{align*}
\beta^*(d_0,d_1,i) & = \fn{rem}(1+(i+1) d_1,d_0) \text{ आणि}\\
\beta(d,i) & = \beta^*(K(d),L(d),i).
\end{align*}
हेच आपल्याला हवे असलेले फलन आहे. वर दिलेल्या
$a_0,\dots,a_n$ साठी
\[j = \max(n,a_0+1,\dots,a_n+1),
\]
असे घ्या, आणि $d_1 = \lcm(1,\dots,j)$ असे घ्या.
वरील \ref{inc:req:bet:rel-prime} नुसार $1+d_1$, $1+2 d_1$,
\dots, $1+(n+1) d_1$ या सापेक्षतः अविभाज्य आहेत,
आणि \ref{inc:req:bet:less} नुसार त्या अनुक्रमे
$a_0,\dots,a_n$ पेक्षा मोठ्या आहेत. सुनझीच्या प्रमेयानुसार
प्रत्येक~$i$ साठी पुढील अट पूर्ण करणारे $d_0$ चे मूल्य आहे:
\[d_0 \equiv a_i \mod (1+(i+1)d_1)
\]
म्हणून ($d_1$ हे~$a_i$ पेक्षा मोठे असल्यामुळे)
\[a_i = \fn{rem}(1+(i+1)d_1,d_0).
\]
$d = J(d_0,d_1)$ असे घ्या. मग प्रत्येक $i \le n$ साठी
\begin{align*}
\beta(d,i) & =  \beta^*(d_0,d_1,i) \\
& =  \fn{rem}(1+(i+1) d_1,d_0) \\
& =  a_i
\end{align*}
हेच अपेक्षित होते. याने $\beta$-फलनावरील साहाय्यक विधानाचा
पुरावा पूर्ण होतो.

\begin{prob}
  $x < y$, $x \mid y$ हे संबंध आणि $\fn{rem}(x,y)$ हे
  फलन आदिम पुनरावर्तनाशिवाय परिभाषित करता येते, हे दाखवा.
  $0$, उत्तरवर्ती फलन, बेरीज, गुणाकार, $\Char{=}$,
  प्रक्षेपणे, तसेच परिबद्ध लघुतमीकरण आणि परिबद्ध संख्यकीकरण
  वापरू शकता.
\end{prob}











\section{आदिम पुनरावर्तनाचे अनुकरण}\label{inc:req:pri:sec}

बीटा फलन वापरून नियमित लघुतमीकरणाने आदिम पुनरावर्तनाने केलेल्या
व्याख्येचे ``अनुकरण'' करता येते, हे आता दाखवू. समजा आपल्याकडे
$f(\vec x)$ आणि $g(\vec x, y, z)$ ही फलने आहेत. मग $f$ आणि~$g$
यांपासून आदिम पुनरावर्तनाने परिभाषित केलेले फलन~$h(\vec x,y)$
असे आहे:
\begin{align*}
h(\vec x, 0) & =  f(\vec x) \\
h(\vec x, y+1) & =  g(\vec x, y, h(\vec x, y)).
\end{align*}
मूलभूत फलने आणि त्यांपासून संयोजन व नियमित लघुतमीकरणाने
परिभाषित केलेली फलने (उदाहरणार्थ~$\beta$) वापरून,
$f$ आणि~$g$ यांपासून केवळ संयोजन व नियमित लघुतमीकरणाने
$h$ परिभाषित करता येते, हे दाखवायचे आहे.

\begin{lem}
\label{inc:req:pri:lem:prim-rec}
$h$ हे $f$ आणि $g$ यांपासून आदिम पुनरावर्तनाने परिभाषित
करता येत असेल, तर $f$, $g$ आणि $\Zero$, $\Succ$,
$\Proj{n}{i}$, $\Add$, $\Mult$, $\Char{=}$ ही फलने
वापरून संयोजन व नियमित लघुतमीकरणानेही ते परिभाषित करता येते.
\end{lem}

\begin{proof}
प्रथम $\hat h(\vec x, y)$ हे साहाय्यक फलन परिभाषित करू.
ते लघुतम संख्या~$d$ परत करते; हा $d$ पुढील अटी
पूर्ण करणाऱ्या क्रमिकेचा संकेतांक असतो:
\begin{enumerate}
\item $(d)_0 = f(\vec x)$, आणि
\item प्रत्येक $i < y$ साठी $(d)_{i+1} = g(\vec x, i, (d)_i)$,
\end{enumerate}
येथे $(d)_i$ हा आता $\beta(d,i)$ चा संक्षेप आहे. दुसऱ्या
शब्दांत, $\hat h$ हे $\tuple{h(\vec x, 0), h(\vec x, 1),
\dots, h(\vec x, y)}$ या क्रमिकेचा संकेतांक परत करते.
$\hat h$ पुढीलप्रमाणे लिहिता येते:
\[\hat h(\vec x, y) = \umin{d}{(\beta(d,0) = f(\vec x) \land \bforall{i <
  y}{\beta(d,i+1) = g(\vec x, i,\beta(d,i)})}.
\]
लक्षात घ्या: येथे आदिम पुनरावर्तन लागत नाही; फक्त
लघुतमीकरण पुरेसे आहे. बीटा फलनावरील साहाय्यक विधान
\ref{inc:req:bet:lem:beta} नुसार या लघुतमीकरणातील शोध नियमित आहे.

आता आपल्याकडे
\[h(\vec x, y) = \beta(\hat h(\vec x, y), y),
\]
म्हणून मूळ दिलेली दोन फलने आणि मूलभूत फलने यांपासून केवळ संयोजन
व नियमित लघुतमीकरणाने $h$ परिभाषित करता येते.
\end{proof}











\section{मूलभूत फलने~$\Th{Q}$ मध्ये निरूपणीय आहेत}\label{inc:req:bre:sec}

सर्वप्रथम सर्व मूलभूत फलने~$\Th{Q}$ मध्ये निरूपणीय आहेत, हे
दाखवायचे आहे. अखेरीस, $k$ आदाने असलेल्या प्रत्येक मूलभूत
फलनासाठी $f(x_0,\dots,x_{k-1})$ त्याचे निरूपण करणारे
सूत्र $\varphi_f(x_0,\dots,x_{k-1},y)$ कसे द्यायचे,
हे दाखवावे लागेल.

शून्य, उत्तरवर्ती, बेरीज, गुणाकार, समतेचे जातिबोधक फल आणि
प्रक्षेपण फलने यांची निरूपणीयता दाखवता येईल. प्रत्येक बाबतीत
योग्य निरूपक सूत्र सहज सुचते; उदाहरणार्थ, शून्याचे निरूपण
$y = \Obj 0$ या सूत्राने, उत्तरवर्ती फलनाचे सूत्र
$x_0' = y$ या सूत्राने आणि बेरीजेचे सूत्र
$(x_0 + x_1) = y$ या सूत्राने होते. खरे काम म्हणजे संबंधित
वाक्ये $\Th{Q}$ मध्ये सिद्ध होतात, हे दाखवणे. उदाहरणार्थ,
वरील सूत्र बेरीजेचे निरूपण करते असे म्हणण्यासाठी
नैसर्गिक संख्यांच्या प्रत्येक जोडी $m$ आणि $n$ साठी $\Th{Q}$
पुढील विधाने सिद्ध करते, हे दाखवावे लागते:
\begin{align*}
& \eq[\num n + \num m][\num {n+m}] \text{ आणि}\\
& \lforall[y][(\eq[(\num n + \num m)][y] \lif \eq[y][\num{n+m}])].
\end{align*}

\begin{prop}
\label{inc:req:bre:prop:rep-zero}
शून्य फलन $\Zero(x) = 0$ चे~$\Th{Q}$ मध्ये
$\varphi_{\Zero}(x,y) \ident \eq[y][\Obj 0]$ या सूत्राने निरूपण होते.
\end{prop}

\begin{prop}
\label{inc:req:bre:prop:rep-succ}
उत्तरवर्ती फलन $\Succ(x) = x+1$ चे~$\Th{Q}$ मध्ये
$\varphi_{\Succ}(x,y) \ident \eq[y][x']$ या सूत्राने निरूपण होते.
\end{prop}

\begin{prop}
\label{inc:req:bre:prop:rep-proj}
प्रक्षेपण फलन $\Proj{n}{i}(x_0, \dots, x_{n-1}) = x_i$ चे
~$\Th{Q}$ मध्ये पुढील सूत्राने निरूपण होते:
\[\varphi_{\Proj{n}{i}}(x_0, \dots, x_{n-1}, y) \ident \eq[y][x_i].\]
\end{prop}

\begin{prob}
$\eq[y][\Obj 0]$, $\eq[y][x']$ आणि $\eq[y][x_i]$
ही सूत्रे अनुक्रमे $\Zero$, $\Succ$ आणि $\Proj{n}{i}$ यांचे
निरूपण करतात, हे सिद्ध करा.
\end{prob}

\begin{prop}
\label{inc:req:bre:prop:rep-id}
समता~$=$ हिचे जातिबोधक फल,
\[\Char{=}(x_0, x_1) =
\begin{cases}
  1 & \text{जर } x_0 =x_1\\
  0 & \text{अन्यथा}
\end{cases}
\]
याचे~$\Th{Q}$ मध्ये पुढील सूत्राने निरूपण होते:
\[  \varphi_{\Char{=}}(x_0, x_1, y) \ident (\eq[x_0][x_1] \land \eq[y][\num{1}]) \lor (\eq/[x_0][x_1] \land
\eq[y][\num{0}]).
\]
\end{prop}

या पुराव्यासाठी पुढील साहाय्यक विधान लागते.

\begin{lem}
\label{inc:req:bre:lem:q-proves-neq} $n$ आणि $m$ या नैसर्गिक संख्या असू द्या.
$n \neq m$ असेल, तर $\Th{Q} \Proves \eq/[\num n][\num m]$.
\end{lem}

\begin{proof}
$n$ वर विगमन वापरून, प्रत्येक $m$ साठी $n \neq m$ असेल तर
$Q \Proves \eq/[\num n][\num m]$, हे दाखवू.

आधार टप्प्यात $n = 0$. जर $m$ हे $0$ च्या बरोबर नसेल,
तर एखादी नैसर्गिक संख्या $k$ अशी असते की $m = k +
1$. $\lforall[x][\eq/[0][x']]$ असे सांगणारे स्वयंसिद्धक
आपल्याकडे आहे. संख्यकाच्या एका स्वयंसिद्धकानुसार $x$ च्या जागी
$\num k$ ठेवून $\eq/[0][\num k']$ हा निष्कर्ष मिळतो. पण
$\num k'$ म्हणजेच $\num m$.

विगमन टप्प्यात, $n$ साठी विधान खरे आहे असे धरून $n+1$ चा
विचार करू. $m$ ही कोणतीही नैसर्गिक संख्या असू द्या. दोन शक्यता
आहेत: $m = 0$, किंवा एखाद्या $k$ साठी $m = k+1$. पहिली
शक्यता वरप्रमाणे हाताळता येते. दुसऱ्या शक्यतेत $n+1 \neq
k+1$ असे समजा. मग $n \neq k$. $n$ साठीच्या विगमन गृहीतकाने
$\Th{Q} \Proves \eq/[\num n][\num k]$. आपल्याकडे
$\lforall[x][\lforall[y][\eq[x'][y'] \lif \eq[x][y]]]$
असे सांगणारे स्वयंसिद्धक आहे. संख्यकाच्या एका स्वयंसिद्धकाने
$\eq[\num n'][\num k'] \lif \eq[\num n][\num
  k]$ मिळते. विधानीय तर्कशास्त्र वापरून $\Th{Q}$ मध्ये
$\eq/[\num n][\num k] \lif \eq/[\num n'][\num k']$
हा निष्कर्ष मिळतो. आता अभिव्यंजन-निर्मूलन नियमाने
$\eq/[\num n'][\num k']$ मिळते; आणि $\num k'$ म्हणजेच
$\num m$ असल्याने आपल्याला हेच दाखवायचे होते.
\end{proof}

\begin{explain}
हे साहाय्यक विधान फार मोठी गोष्ट सांगत नाही: वेगवेगळे संख्यांक
वेगवेगळ्या वस्तू दर्शवतात, हे $\Th{Q}$ सिद्ध करू शकते,
असे त्याचे सार आहे. उदाहरणार्थ, $\Th{Q}$ हे
$0'' \neq 0'''$ सिद्ध करते. पण हे सर्वसाधारणपणे खरे आहे
हे दाखवताना काही काळजी घ्यावी लागते. येथे विगमन वापरले असले,
तरी ते $\Th{Q}$ च्या \emph{बाहेर} केलेले विगमन आहे, हेही लक्षात घ्या.
\end{explain}

\begin{proof}[\ref{inc:req:bre:prop:rep-id} चा पुरावा]
$n = m$ असेल, तर $\num{n}$ आणि $\num{m}$ ही एकच पदे आहेत
आणि $\Char{=}(n, m) = 1$. तसेच $\Th{Q} \Proves (\eq[\num{n}][\num{m}] \land
\eq[\num{1}][\num{1}])$, म्हणून ते $\varphi_=(\num{n}, \num{m},
\num{1})$ सिद्ध करते. $n \neq m$ असेल, तर $\Char=(n, m) = 0$.
\ref{inc:req:bre:lem:q-proves-neq} नुसार $\Th{Q} \Proves \eq/[\num{n}][\num{m}]$;
म्हणून $(\eq/[\num{n}][\num{m}] \land \Obj 0 = \Obj 0)$
हेही सिद्ध होते. त्यामुळे $\Th{Q}
\Proves \varphi_=(\num{n}, \num{m}, \num{0})$.

दुसऱ्या अटीसाठीही दोन शक्यता आहेत. $n = m$ असेल, तर
$\Th{Q} \Proves \lforall[y][(\varphi_=(\num{n}, \num{m}, y)
  \lif \eq[y][\num{1}])]$ हे दाखवायचे आहे. अनौपचारिकपणे
युक्तिवाद करू: $\varphi_=(\num{n}, \num{m}, y)$, म्हणजेच
\[(\eq[\num{n}][\num{n}] \land \eq[y][\num{1}]) \lor
(\eq/[\num{n}][\num{n}] \land \eq[y][\num{0}])
\]
असे समजा. डाव्या विकल्पातून तर्कशास्त्राने
$\eq[y][\num{1}]$ मिळते. उजवा विकल्प मात्र
तर्कशास्त्राने सिद्ध होणाऱ्या $\eq[\num{n}][\num{n}]$
शी विसंगत आहे.

दुसरीकडे, $n \neq m$ असेल, तर $\varphi_=(\num{n},
\num{m}, y)$ असे आहे:
\[(\eq[\num{n}][\num{m}] \land \eq[y][\num{1}]) \lor
(\eq/[\num{n}][\num{m}] \land \eq[y][\num{0}])
\]
येथे डावा विकल्प \ref{inc:req:bre:lem:q-proves-neq} नुसार
$\Th{Q}$ मध्ये सिद्ध होणाऱ्या $\eq/[\num{n}][\num{m}]$
शी विसंगत आहे; उजव्या विकल्पातून $\eq[y][\num{0}]$ मिळते.
\end{proof}

\begin{prop}
\label{inc:req:bre:prop:rep-add}
बेरीज फलन $\Add(x_0, x_1) = x_0+x_1$ चे~$\Th{Q}$ मध्ये
पुढील सूत्राने निरूपण होते:
\[  \varphi_{\Add}(x_0, x_1, y) \ident \eq[y][(x_0 + x_1)].
\]
\end{prop}

\begin{lem}
\label{inc:req:bre:lem:q-proves-add}
$\Th{Q} \Proves \eq[(\num{n} + \num{m})][\num{n+m}]$
\end{lem}

\begin{proof}
$m$ वर विगमन वापरून हे सिद्ध करू. $m = 0$ असेल, तर
$\Th{Q} \Proves \eq[(\num{n} + \Obj 0)][\num{n}]$
हे विधान आहे. ते स्वयंसिद्धक~$\mathsf{Q}_4$ वरून मिळते. आता
$m$ साठी विधान खरे आहे असे समजून $m+1$ साठी, म्हणजे
$\Th{Q} \Proves \eq[(\num{n} +
  \num{m+1})][\num{n+m+1}]$ हे सिद्ध करू. लक्षात घ्या,
$\num{m+1}$ म्हणजेच $\num{m}'$ आणि $\num{n+m+1}$ म्हणजेच
$\num{n+m}'$. स्वयंसिद्धक $\mathsf{Q}_5$ ने $\Th{Q}
\Proves \eq[(\num{n} + \num{m}')][(\num{n}+\num{m})']$.
विगमन गृहीतकाने $\Th{Q} \Proves \eq[(\num{n} + \num{m})][\num{n+m}]$.
म्हणून $\Th{Q} \Proves \eq[(\num{n} + \num{m}')][\num{n+m}']$.
\end{proof}

\begin{proof}[\ref{inc:req:bre:prop:rep-add} चा पुरावा]
$\Add$ चे निरूपण करणारे सूत्र~$\varphi_\Add(x_0, x_1, y)$
म्हणजे $\eq[y][(x_0 + x_1)]$. प्रथम, $\Add(n, m) = k$
असेल, तर $\Th{Q} \Proves \varphi_\Add(\num{n}, \num{m}, \num{k})$,
म्हणजे $\Th{Q}
\Proves \eq[\num{k}][(\num{n} + \num{m})]$, हे दाखवू.
पण $k = n + m$ असल्यामुळे $\num{k}$ म्हणजेच $\num{n+m}$;
आणि \ref{inc:req:bre:lem:q-proves-add} मध्ये आपण
$\Th{Q} \Proves \eq[(\num{n} +
  \num{m})][\num{n+m}]$ हे सिद्ध केले आहे. समतेच्या
सममितीने त्यावरून अपेक्षित उलट क्रमातील समता मिळते.

$\Add(n, m) = k$ असेल, तर पुढील विधानही दाखवायचे आहे:
\[\Th{Q} \Proves \lforall[y][(\varphi_\Add(\num{n}, \num{m}, y) \lif
  \eq[y][\num{k}])].
\]
निरूपक सूत्र गृहीत धरून, समतेच्या सममितीने
$\eq[(\num{n} + \num{m})][y]$ असे समजा. कारण
\[\Th{Q} \Proves \eq[(\num{n}+\num{m})][\num{n+m}],
\]
म्हणून डावीकडील पदाऐवजी $\num{n+m}$ ठेवून कोणत्याही~$y$
साठी $\eq[\num{n+m}][y]$ मिळते. समतेच्या सममितीने
अपेक्षित क्रमातील समता मिळते; उलट अभिव्यंजनही याच सिद्ध
समतेवरून आणि समतेच्या नियमांवरून मिळते.
\end{proof}

\begin{prop}
\label{inc:req:bre:prop:rep-mult}
गुणाकार फलन $\Mult(x_0, x_1) = x_0 \cdot x_1$ चे
~$\Th{Q}$ मध्ये पुढील सूत्राने निरूपण होते:
\[  \varphi_{\Mult}(x_0, x_1, y) \ident y = (x_0 \times x_1).
\]
\end{prop}

\begin{proof} सरावासाठी. \end{proof}

\begin{lem}
\label{inc:req:bre:lem:q-proves-mult}
$\Th{Q} \Proves \eq[(\num{n} \times \num{m})][\num{n \cdot m}]$
\end{lem}

\begin{proof} सरावासाठी. \end{proof}

\begin{prob}
\ref{inc:req:bre:lem:q-proves-mult} सिद्ध करा.
\end{prob}

\begin{prob}
\ref{inc:req:bre:lem:q-proves-mult} वापरून
\ref{inc:req:bre:prop:rep-mult} सिद्ध करा.
\end{prob}

\begin{explain}
  आठवून पाहा: अंकगणिताच्या भाषेतील फलनचिन्हासाठी आपण $\times$
  वापरतो, तर संख्यांवरील नेहमीच्या गुणाकारासाठी $\cdot$ वापरतो.
  म्हणून संख्यादर्शक व्यंजकांमध्ये (उदाहरणार्थ, $m \cdot n$)
  $\cdot$ येऊ शकते; पण अंकगणिताच्या भाषेतील पदांमध्येच
  (उदाहरणार्थ, $(\num{m} \times
  \num{n})$) $\times$ येते. आणखी गोंधळाची बाब म्हणजे $+$
  हे चिन्ह फलन आणि बेरीजक्रिया या दोन्हींसाठी वापरतो.
  ते पदांच्या मध्ये आल्यास---उदाहरणार्थ, $(\num{n} + \num{m})$
  मध्ये---अंकगणिताच्या भाषेतील $2$ आदानांचे फलन असते;
  आणि संख्यांच्या मध्ये आल्यास---उदाहरणार्थ, $n+m$ मध्ये---
  ती बेरीजक्रिया असते. $\num{n+m}$ हे उदाहरणही यांत येते:
  तो $n+m$ या संख्येला अनुरूप असलेला मानक संख्यांक आहे.
\end{explain}











\section{फलनांचे संयोजन~$\Th{Q}$ मध्ये निरूपणीय असते}\label{inc:req:cmp:sec}

समजा $h$ पुढीलप्रमाणे परिभाषित आहे:
\[h(x_0,\dots,x_{l-1}) = f(g_0(x_0,\dots,x_{l-1}), \dots,
g_{k-1}(x_0,\dots,x_{l-1})).
\]
येथे $f$ आणि अनुक्रमे $g_0$,
\dots,~$g_{k-1}$ या फलनांचे निरूपण करणारी
सूत्रे $\varphi_f, \varphi_{g_0}, \dots,
\varphi_{g_{k-1}}$ आपल्याला आधीच सापडली आहेत. आता $h$ चे
निरूपण करणारे सूत्र~$\varphi_h$ शोधायचे आहे.

प्रत्येक फलन $1$-स्थानी असलेली सोपी बाब आधी पाहू; म्हणजे
$h(x) = f(g(x))$ विचारात घेऊ. $\varphi_f(y, z)$ हे~$f$ चे,
आणि $\varphi_g(x, y)$ हे~$g$ चे निरूपण करत असेल, तर $h$ चे
निरूपण करणारे सूत्र~$\varphi_h(x, z)$ शोधायचे आहे.
$h(x) = z$, तर आणि तरच असा एखादा~$y$ असतो की
$z = f(y)$ आणि $y = g(x)$ ही दोन्ही विधाने खरी असतात.
($h(x) = z$ असेल तर $g(x)$ हा असा~$y$ असतो; आणि असा
$y$ असेल, तर $y = g(x)$ आणि $z = f(y)$ असल्याने
$z = f(g(x))$.) यावरून $\varphi_h(x, z)$ साठी
$\lexists[y][(\varphi_g(x, y) \land \varphi_f(y,
  z))]$ हे योग्य सूत्र असावे असे दिसते. आता फक्त
$\Th{Q}$ संबंधित सूत्रे सिद्ध करते, हे तपासायचे आहे.

\begin{prop}
\label{inc:req:cmp:prop:rep1}
$h(n) = m$ असेल, तर $\Th{Q} \Proves \varphi_h(\num{n}, \num{m})$.
\end{prop}

\begin{proof}
$h(n) = m$, म्हणजे $f(g(n)) = m$, असे समजा. $k = g(n)$ असू द्या.
मग
\begin{align*}
  \Th{Q} & \Proves \varphi_g(\num{n}, \num{k})
  \intertext{कारण $\varphi_g$ हे~$g$ चे निरूपण करते; तसेच}
  \Th{Q} & \Proves \varphi_f(\num{k}, \num{m})
  \intertext{कारण $\varphi_f$ हे~$f$ चे निरूपण करते. म्हणून}
  \Th{Q} & \Proves \varphi_g(\num{n}, \num{k}) \land \varphi_f(\num{k}, \num{m})
  \intertext{आणि त्यावरून}
  \Th{Q} & \Proves \lexists[y][(\varphi_g(\num{n}, y) \land \varphi_f(y, \num{m}))],
\end{align*}
म्हणजेच $\Th{Q} \Proves \varphi_h(\num{n}, \num{m})$.
\end{proof}

\begin{prop}
\label{inc:req:cmp:prop:rep2}
$h(n) = m$ असेल, तर $\Th{Q} \Proves \lforall[z][(\varphi_h(\num{n}, z) \lif
  z = \num{m})]$.
\end{prop}

\begin{proof}
$h(n) = m$, म्हणजे $f(g(n)) = m$, असे समजा. $k = g(n)$ असू द्या.
मग
\begin{align*}
  \Th{Q} & \Proves \lforall[y][(\varphi_g(\num{n}, y) \lif \eq[y][\num{k}])]
  \intertext{कारण $\varphi_g$ हे~$g$ चे निरूपण करते; तसेच}
  \Th{Q} & \Proves \lforall[z][(\varphi_f(\num{k}, z) \lif \eq[z][\num{m}])]
  \intertext{कारण $\varphi_f$ हे~$f$ चे निरूपण करते. थोडेसे
    तर्कशास्त्र वापरून पुढील विधानही दाखवता येते:}
  \Th{Q} & \Proves \lforall[z][(\lexists[y][(\varphi_g(\num{n}, y) \land
      \varphi_f(y, z))] \lif \eq[z][\num{m}])].
\end{align*}
म्हणजेच $\Th{Q} \Proves \lforall[y][(\varphi_h(\num n, y) \lif \eq[y][\num m])]$.
\end{proof}

$f$ आणि~$g_i$ यांची आदाने $1$ पेक्षा अधिक असलेल्या गुंतागुंतीच्या
बाबतीतही हीच कल्पना लागू पडते.

\begin{prop}
\label{inc:req:cmp:prop:rep-composition}
$\varphi_f(y_0, \dots, y_{k-1}, z)$ हे~$\Th{Q}$ मध्ये
$f(y_0, \dots, y_{k-1})$ चे, आणि
$\varphi_{g_i}(x_0, \dots, x_{l-1}, y)$ हे~$\Th{Q}$ मध्ये
$g_i(x_0, \dots, x_{l-1})$ चे निरूपण करत असेल, तर
\begin{multline*}
  \lexists[y_0\dots][\lexists[y_{k-1}][(\varphi_{g_0}(x_0,\dots,x_{l-1},y_0) \land
      \dots \land {}]]\\
  \varphi_{g_{k-1}}(x_0,\dots,x_{l-1},y_{k-1}) \land \varphi_f(y_0,\dots,y_{k-1},z))
\end{multline*}
हे सूत्र पुढील फलनाचे निरूपण करते:
\[h(x_0, \dots, x_{l-1}) = f(g_0(x_0, \dots, x_{l-1}), \dots, g_{k-1}(x_0,
\dots, x_{l-1})).
\]
\end{prop}

\begin{proof}
सरावासाठी.
\end{proof}

\begin{prob}
\ref{inc:req:cmp:prop:rep1} आणि
\ref{inc:req:cmp:prop:rep2} यांच्या पुराव्यांचा आधार घेऊन
\ref{inc:req:cmp:prop:rep-composition} चा पुरावा तपशीलवार द्या.
\end{prob}











\section{नियमित लघुतमीकरण~$\Th{Q}$ मध्ये निरूपणीय असते}\label{inc:req:min:sec}

अपरिबद्ध शोधाचा विचार करू. समजा $g(x, z)$ हे नियमित फलन
$\Th{Q}$ मध्ये, उदाहरणार्थ सूत्र~$\varphi_g(x, z, y)$ या
सूत्राने, निरूपणीय आहे. $f(z) = \umin{x}{[g(x, z) = 0]}$
असे $f$ परिभाषित करू. $f$ चे निरूपण करणारे
सूत्र~$\varphi_f(z, y)$ शोधायचे आहे. $f(z)$ चे मूल्य
म्हणजे अशी संख्या~$x$ की (अ) $g(x, z) = 0$ आणि (आ) तशी
ती लघुतम आहे, म्हणजे कोणत्याही $w < x$ साठी
$g(w, z) \neq 0$. म्हणून पुढील निवड स्वाभाविक आहे:
\[\varphi_f(z,y) \ident \varphi_g(y, z, \Obj 0) \land \lforall[w][(w < y \lif \lnot
  \varphi_g(w, z, \Obj 0))].
\]
अर्थात, सर्वसाधारण बाबतीत $z$ ऐवजी $z_0$, \dots, $z_k$
ही चलं घ्यावी लागतील.

या पुराव्यातही $\Th{Q}$ ची सिद्ध करण्याची शक्ती किती आहे,
याविषयी काही साहाय्यक विधाने लागतील.

\begin{lem}
\label{inc:req:min:lem:succ} प्रत्येक स्थिरांक~$a$ आणि प्रत्येक नैसर्गिक
संख्या~$n$ साठी,
\[\Th{Q} \Proves \eq[(a' + \num n)][(a + \num n)'].
\]
\end{lem}

\begin{proof}
नेहमीप्रमाणे $n$ वर विगमन वापरू. आधार टप्प्यात $n =
0$ असेल, तर $\Th{Q}$ हे $\eq[(a' + \Obj 0)][(a + \Obj
0)']$ सिद्ध करते, हे दाखवायचे आहे. पण पुढील विधाने मिळतात:
\begin{align}
  \Th{Q} & \Proves \eq[(a' + \Obj 0)][a'] \quad \text{स्वयंसिद्धक $Q_4$ ने}
  \label{inc:req:min:step1}\\
  \Th{Q} & \Proves  \eq[(a + \Obj 0)][a] \quad \text{स्वयंसिद्धक $Q_4$ ने}
  \label{inc:req:min:step2} \\
  \Th{Q} & \Proves \eq[(a + \Obj 0)'][a'] \quad \text{\ref{inc:req:min:step2} वरून}
  \label{inc:req:min:step3} \\
  \Th{Q} & \Proves \eq[(a' + \Obj 0)][(a + \Obj 0)'] \quad
  \text{\ref{inc:req:min:step1} आणि \ref{inc:req:min:step3} वरून}\notag
\end{align}
विगमन टप्प्यात, $\Th{Q}
\Proves \eq[(a' + \num n)][(a + \num n)']$ हे सिद्ध झाले आहे
असे समजू. $\num{n+1}$ म्हणजे $\num{n}'$ असल्याने, आता
$\Th{Q}$ हे $\eq[(a' + \num{n}')][(a + \num{n}')']$
सिद्ध करते, हे दाखवायचे आहे. आपल्याला मिळते:
\begin{align}
  \Th{Q} & \Proves \eq[(a' + \num n')][(a' + \num n)'] \quad
  \text{स्वयंसिद्धक $\mathsf{Q}_5$ ने} \label{inc:req:min:step5}\\
  \Th{Q} & \Proves \eq[(a' + \num n)'][(a + \num n')'] \quad
  \text{विगमन गृहीतक, बेरीजेचे स्वयंसिद्धक आणि समतेचे नियम वापरून} \label{inc:req:min:step6}\\
  \Th{Q} & \Proves \eq[(a' + \num n')][(a + \num n')'] \quad
  \text{\ref{inc:req:min:step5} आणि \ref{inc:req:min:step6} वरून.} \notag
\end{align}
\end{proof}

हे पुन्हा लक्षात घ्यावे की यावरून $\Th{Q}$ हे
$\lforall[x][\lforall[y][(x' + y) = (x + y)']]$ सिद्ध करते,
असे म्हणता येत नाही; वरचे विधान त्याहून दुर्बल आहे. हे
वाक्य~$\Struct{N}$ मध्ये खरे असले, तरी $\Th{Q}$
त्याला सिद्ध करू शकत नाही.

\begin{lem}
\label{inc:req:min:lem:less-zero}
$\Th{Q} \Proves \lforall[x][\lnot x < \Obj 0]$.
\end{lem}

\begin{proof}
पुरावा अनौपचारिक रीतीने देऊ (म्हणजे औपचारिक निष्पत्ती
कशी रचावी यासाठी फक्त संकेत देऊ).

कोणत्याही~$a$ साठी $\lnot a < \Obj 0$ हे सिद्ध करायचे आहे.
$<$ च्या व्याख्येनुसार $\Th{Q}$ मध्ये
$\lnot \lexists[y][\eq[(y'+a)][\Obj 0]]$ सिद्ध करायचे आहे.
$\lexists[y][\eq[(y'+a)][\Obj 0]]$ असे मानून व्याघात मिळवू.
$\eq[(b'+a)][\Obj 0]$ असे समजा. $\mathsf{Q}_3$ वापरल्यावर
$\eq[a][\Obj 0] \lor \lexists[y][\eq[a][y']]$ मिळते.
आता दोन बाबी पाहू.

पहिली बाब: $\eq[a][\Obj 0]$ खरे आहे. $\eq[(b'+a)][\Obj 0]$
वरून $\eq[(b' + \Obj 0)][\Obj 0]$ मिळते. $\Th{Q}$ च्या
स्वयंसिद्धक~$\mathsf{Q}_4$ ने $\eq[(b' + \Obj 0)][b']$ मिळते,
म्हणून $\eq[b'][\Obj 0]$ मिळते. पण स्वयंसिद्धक~$\mathsf{Q}_2$ ने
$\eq/[b'][\Obj 0]$ हेसुद्धा मिळते; हा व्याघात आहे.

दुसरी बाब: एखाद्या $c$ साठी $\eq[a][c']$. मग
$\eq[(b' +
c')][\Obj 0]$ मिळते. स्वयंसिद्धक~$\mathsf{Q}_5$ ने
$\eq[(b' + c)'][\Obj 0]$ मिळते; हेसुद्धा स्वयंसिद्धक~$Q_2$
शी विसंगत आहे.
\end{proof}

\begin{lem}
\label{inc:req:min:lem:less-nsucc}
प्रत्येक नैसर्गिक संख्या~$n$ साठी,
  \[  \Th{Q} \Proves
  \lforall[x][(x < \num {n+1} \lif (\eq[x][\Obj 0] \lor \dots \lor
    \eq[x][\num n]))].
  \]
\end{lem}

\begin{proof}
$n$ वर विगमन वापरू. प्रथम $n = 0$ हा आधार टप्पा पाहू.
त्या वेळी कोणत्याही~$a$ साठी $a < \num 1 \lif \eq[a][\Obj 0]$
हे दाखवायचे आहे. $a < \num 1$ असे समजा. मग $<$ च्या
परिभाषक स्वयंसिद्धकाने $\lexists[y][\eq[(y'+a)][\Obj 0']]$
मिळते (कारण $\num 1 \ident
\Obj 0'$).

$b$ कडे हा गुणधर्म आहे, म्हणजे $\eq[(b'+a)][\Obj 0']$,
असे समजा. $\eq[a][\Obj 0]$ दाखवायचे आहे. स्वयंसिद्धक~$\mathsf{Q}_3$
ने $\eq[a][\Obj 0]$ किंवा एखाद्या $c$ साठी
$\eq[a][c']$ मिळते. पहिल्या बाबतीत आणखी काही दाखवायचे नाही.
म्हणून $\eq[a][c']$ असे समजा. मग
$\eq[(b' + c')][\Obj 0']$ मिळते. $\Th{Q}$ च्या स्वयंसिद्धक~$\mathsf{Q}_5$
ने $\eq[(b'+c)'][\Obj 0']$ मिळते. स्वयंसिद्धक~$\mathsf{Q}_1$ ने
$\eq[(b' + c)][\Obj
0]$ मिळते. पण स्वयंसिद्धक~$\mathsf{Q}_8$ नुसार याचा अर्थ
$c < \Obj 0$; याला \ref{inc:req:min:lem:less-zero} ने विरोध होतो.

आता विगमन टप्पा पाहू. $n$ साठीचे विधान गृहीत धरून $n+1$
साठी ते सिद्ध करू. $a < \num {n+2}$ असे समजा. पुन्हा
$\mathsf{Q}_3$ वापरून दोन बाबी मिळतात: $\eq[a][\Obj 0]$, किंवा
एखाद्या $b$ साठी $\eq[a][b']$. पहिल्या बाबतीत
$\eq[a][\Obj 0] \lor \dots \lor
\eq[a][\num{n+1}]$ सहज मिळते. दुसऱ्या बाबतीत $b'
< \num {n+2}$, म्हणजे $b' < \num{n+1}'$. स्वयंसिद्धक~$\mathsf{Q}_8$
नुसार एखाद्या $c$ साठी
$\eq[(c'+b')][\num{n+1}']$. स्वयंसिद्धक $\mathsf{Q}_5$ ने
$\eq[(c'+b)'][\num{n+1}']$. स्वयंसिद्धक~$\mathsf{Q}_1$ ने
$\eq[(c'+b)][\num{n+1}]$; म्हणून स्वयंसिद्धक~$\mathsf{Q}_8$ ने
$b < \num{n+1}$. विगमन गृहीतकाने
$\eq[b][\Obj 0] \lor \dots \lor \eq[b][\num{n}]$.
यावरून तर्कशास्त्राने
$\eq[b'][\Obj 0'] \lor \dots \lor \eq[b'][\num{n}']$
मिळते; आणि $\eq[a][b']$ असल्याने
$\eq[a][\num{1}] \lor \dots \lor \eq[a][\num{n+1}]$
मिळते. उलट दिशेसाठी, उजवीकडील प्रत्येक निर्दिष्ट संख्यांकाशी
समता गृहीत धरल्यास, वरच्या मर्यादेपर्यंतचे उरलेले अंतर
दाखवणारी विशिष्ट संख्या साक्षी म्हणून घ्यावी. आधी सिद्ध
केलेल्या संख्यांकांच्या बेरीजविषयक विधानाने अपेक्षित बेरीज
सिद्ध होते आणि क्रमसंबंधाच्या परिभाषक स्वयंसिद्धकाने तो संख्यांक
वरच्या मर्यादेपेक्षा लहान ठरतो;
आधार टप्प्यातही याच रीतीने उलट दिशा मिळते.
\end{proof}

\begin{lem}
  \label{inc:req:min:lem:trichotomy} प्रत्येक नैसर्गिक संख्या~$m$ साठी,
  \[  \Th{Q} \Proves
  \lforall[y][((y < \num{m} \lor \num{m} < y) \lor \eq[y][\num{m}])].
  \]
\end{lem}

\begin{proof}
$m$ वर विगमन वापरू. प्रथम $m=0$ ही बाब पाहू. $\mathsf{Q}_3$ ने
$\Th{Q} \Proves
\lforall[y][(\eq[y][\Obj 0] \lor \lexists[z][\eq[y][z']])]$.
$a$ कोणताही असू द्या. मग $\eq[a][\Obj 0]$ किंवा एखाद्या~$b$
साठी $\eq[a][b']$. पहिल्या बाबतीत
$(a < \Obj 0 \lor \Obj
0 < a) \lor \eq[a][\Obj 0]$ हेसुद्धा मिळते. पण
$\eq[a][b']$ असेल, तर $\eq$ च्या तर्कशास्त्राने
$\eq[(b' +
\Obj 0)][(a + \Obj 0)]$ मिळते. $\mathsf{Q}_4$ ने
$\eq[(a +
\Obj 0)][a]$, म्हणून $\eq[(b' + \Obj 0)][a]$ आणि त्यामुळे
$\lexists[z][\eq[(z' + \Obj 0)][a]]$ मिळते. $\mathsf{Q}_8$
मधील $<$ च्या व्याख्येनुसार $\Obj 0 < a$. $\Obj 0 < a$
असेल, तर $(\Obj 0 < a \lor a
< \Obj 0) \lor \eq[a][\Obj 0]$ हेसुद्धा मिळते.

आता पुढील विधान गृहीत धरू:
\begin{align*}
  \Th{Q} & \Proves \lforall[y][((y < \num{m} \lor \num{m} < y) \lor
    \eq[y][\num{m}])]
  \intertext{आणि पुढील विधान सिद्ध करायचे आहे:}
  \Th{Q} & \Proves \lforall[y][((y < \num{m+1} \lor \num{m+1} < y) \lor
    \eq[y][\num{m+1}])]
\end{align*}
$a$ कोणताही असू द्या. $\mathsf{Q}_3$ ने $\eq[a][\Obj 0]$ किंवा
एखाद्या~$b$ साठी $\eq[a][b']$. पहिल्या बाबतीत $\mathsf{Q}_4$ ने
$\eq[\num{m}' +
a][\num{m+1}]$ मिळते; म्हणून $\mathsf{Q}_8$ ने
$a < \num{m+1}$.

आता दुसरी बाब, $\eq[a][b']$, पाहू. विगमन गृहीतकानुसार
$(b < \num{m} \lor \num{m} < b) \lor
    \eq[b][\num{m}]$.

पहिला विकल्प $b < \num{m}$ हा $\mathsf{Q}_8$ नुसार
$\lexists[z][\eq[(z' + b)][\num{m}]]$ याच्याशी समतुल्य आहे.
एखाद्या $c$ कडे हा गुणधर्म आहे असे समजा.
$\eq[(c' + b)][\num{m}]$ असेल, तर
$\eq[(c' + b)'][\num{m}']$ हेसुद्धा मिळते. $\mathsf{Q}_5$ ने
$\eq[(c' + b)'][(c' + b')]$ मिळते. म्हणून
$\eq[(c'+b')][\num{m}']$ मिळते. $\num{m}'
\ident \num{m+1}$ हे लक्षात ठेवून $c$ वर अस्तित्ववाचक
सामान्यीकरण केल्यावर
$\lexists[u][\eq[(u' + b')][\num{m+1}]]$ मिळते.
म्हणून $b < \num{m}$ असेल, तर $b' < \num{m+1}$, आणि
त्यामुळे $a < \num{m+1}$.

आता $\num{m} < b$ असे समजा, म्हणजे
$\lexists[z][\eq[(z' +
\num{m})][b]]$. $c$ हा असा~$z$ असेल, म्हणजे
$\eq[(c' +
\num{m})][b]$, असे समजा. तर्कशास्त्राने
$\eq[(c' + \num{m})'][b']$ मिळते. $\mathsf{Q}_5$ ने
$\eq[(c' + \num{m}')][b']$ मिळते. $\eq[a][b']$ आणि
$\num{m}' \ident
\num{m+1}$ असल्याने $\eq[(c' + \num{m+1})][a]$.
$\mathsf{Q}_8$ ने $\num{m+1} < a$.

शेवटी $\eq[b][\num{m}]$ असे समजू. मग तर्कशास्त्राने
$\eq[b'][\num{m}']$, आणि म्हणून
$\eq[a][\num{m+1}]$ मिळते.

अशा प्रकारे $m$ आणि~$b$ यांच्या बाबतीतील प्रत्येक विकल्पापासून
$m+1$ आणि~$a$ यांच्या बाबतीतील अनुरूप विकल्प मिळतो.
\end{proof}

\begin{prop}
\label{inc:req:min:prop:rep-minimization}
$\varphi_g(x, z, y)$ हे $g(x, z)$ चे~$\Th{Q}$ मध्ये निरूपण
करत असेल आणि ते फलन नियमित असेल, तर
\[\varphi_f(z,y) \ident \varphi_g(y, z, \Obj 0) \land \lforall[w][(w < y \lif \lnot
  \varphi_g(w, z, \Obj 0))]
\]
हे सूत्र $f(z) = \umin{x}{[g(x, z) = 0]}$ चे निरूपण करते.
\end{prop}

\begin{proof}
प्रथम $f(n) = m$ असेल, तर
$\Th{Q} \Proves \varphi_f(\num n, \num m)$, म्हणजे पुढील विधान,
हे दाखवू:
\begin{align}
  \Th{Q} & \Proves \varphi_g(\num{m}, \num{n}, \Obj 0) \land \lforall[w][(w <
    \num{m} \lif \lnot \varphi_g(w, \num{n}, \Obj 0))]. \notag
  \intertext{$\varphi_g(x, z, y)$ हे $g(x, z)$ चे निरूपण करते आणि
    $f(n) = m$ असेल तेव्हा $g(m, n) =
    0$; म्हणून}
\Th{Q} & \Proves \varphi_g(\num{m}, \num{n}, \Obj 0). \notag
\intertext{$f(n) = m$ असेल, तर प्रत्येक $k < m$ साठी
  $g(k, n) \neq 0$. निरूपणातील एकमेव-मूल्याची अट आणि
  भिन्न संख्यांकांची असमता वापरल्यावर}
\Th{Q} & \Proves \lnot \varphi_g(\num{k}, \num{n}, \Obj 0). \notag
\intertext{यावरून पुढील विधान मिळते:}
\Th{Q} & \Proves \lforall[w][(w < \num{m} \lif \lnot
  \varphi_g(w, \num{n}, \Obj 0))]. \label{inc:req:min:rep-less}
\end{align}
$m = 0$ असेल, तर हे \ref{inc:req:min:lem:less-zero} वरून;
अन्यथा \ref{inc:req:min:lem:less-nsucc} वरून मिळते.

आता $f(n) = m$ असेल, तर
$\Th{Q} \Proves
\lforall[y][(\varphi_f(\num{n}, y) \lif \eq[y][\num{m}])]$
हे दाखवू. औपचारिक सिद्धता वाचकावर सोडून युक्तिवाद पुन्हा
अनौपचारिकपणे सांगू.

$\varphi_f(\num{n}, b)$ असे समजा. त्यावरून (अ)
$\varphi_g(b, \num{n},
\Obj 0)$ आणि (आ) $\lforall[w][(w < b \lif \lnot \varphi_g(w, \num{n}, \Obj
  0))]$ मिळते. \ref{inc:req:min:lem:trichotomy} नुसार
$(b < \num{m} \lor \num{m} < b)
\lor \eq[b][\num{m}]$. $b < \num{m}$ किंवा $\num{m}
< b$ या दोन्ही शक्यतांपासून व्याघात मिळतो, हे दाखवू.

$\num{m} < b$ असेल, तर (आ) वरून
$\lnot \varphi_g(\num{m}, \num{n}, \Obj 0)$ मिळते. पण
$m = f(n)$ असल्याने $g(m, n) = 0$; आणि $\varphi_g$ हे~$g$
चे निरूपण करत असल्याने $\Th{Q} \Proves
\varphi_g(\num{m}, \num{n}, \Obj 0)$. हा व्याघात आहे.

आता $b < \num{m}$ असे समजा. \ref{inc:req:min:rep-less} नुसार
$\Th{Q} \Proves \lforall[w][(w <
  \num{m} \lif \lnot \varphi_g(w, \num{n}, \Obj 0))]$,
म्हणून $\lnot \varphi_g(b, \num{n}, \Obj 0)$ मिळते.
हे पुन्हा (अ) शी विसंगत आहे. उलट अभिव्यंजन आधी सिद्ध
केलेल्या निरूपक सूत्राच्या निर्दिष्ट मूल्याच्या उदाहरणावरून आणि
समतेच्या प्रतिस्थापन नियमावरून मिळते.
\end{proof}











\section{संगणनक्षम फलने~$\Th{Q}$ मध्ये निरूपणीय असतात}\label{inc:req:crq:sec}

\begin{thm}
प्रत्येक संगणनक्षम फलन~$\Th{Q}$ मध्ये निरूपणीय असते.
\end{thm}

\begin{proof}
संदिग्धता टाळण्यासाठी, चर्च--ट्यूरिंग प्रबंध वापरून फलन
संगणनक्षम असते तर आणि तरच ते सामान्य पुनरावर्ती असते,
असे म्हणू. सामान्य पुनरावर्ती फलने म्हणजे शून्य फलन~$\Zero$,
उत्तरवर्ती फलन~$\Succ$ आणि प्रक्षेपण फलन~$\Proj{n}{i}$
यांपासून संयोजन, आदिम पुनरावर्तन व नियमित लघुतमीकरण वापरून
परिभाषित करता येणारी फलने. \ref{inc:req:pri:lem:prim-rec} नुसार
$f$ आणि~$g$ यांपासून आदिम पुनरावर्तनाने परिभाषित होणारे
कोणतेही फलन~$h$ हे $f$, $g$, $\Zero$,
$\Succ$, $\Proj{n}{i}$, $\Add$, $\Mult$, $\Char{=}$
यांपासून संयोजन व नियमित लघुतमीकरण वापरूनही परिभाषित
करता येते. परिणामी, फलन सामान्य पुनरावर्ती असते तर आणि तरच ते $\Zero$, $\Succ$, $\Proj{n}{i}$, $\Add$,
$\Mult$, $\Char{=}$ यांपासून संयोजन व नियमित लघुतमीकरण
वापरून परिभाषित करता येते.

शिवाय, ही सर्व मूलभूत फलने~$\Th{Q}$ मध्ये निरूपणीय आहेत,
हे आपण सिद्ध केले आहे
(\ref{inc:req:bre:prop:rep-zero} आणि \ref{inc:req:bre:prop:rep-succ} आणि \ref{inc:req:bre:prop:rep-proj} आणि \ref{inc:req:bre:prop:rep-id} आणि \ref{inc:req:bre:prop:rep-add} आणि \ref{inc:req:bre:prop:rep-mult});
आणि निरूपणीय फलनांपासून संयोजन किंवा नियमित लघुतमीकरणाने
परिभाषित होणारे कोणतेही फलन
(\ref{inc:req:cmp:prop:rep-composition},
\ref{inc:req:min:prop:rep-minimization}) हेही निरूपणीय असते.
म्हणून प्रत्येक सामान्य पुनरावर्ती फलन~$\Th{Q}$ मध्ये
निरूपणीय असते.
\end{proof}

\begin{explain}
संगणनक्षम फलनांचा वर्ग म्हणजेच $\Th{Q}$ मध्ये निरूपणीय
फलनांचा वर्ग, असे लक्षणन आपल्याला मिळाले आहे. प्रत्यक्षात
हा पुरावा अधिक व्यापक आहे. निरूपणीयतेच्या व्याख्येवरून
$\Th{Q}$ चा विस्तार असलेल्या कोणत्याही उपपत्तीत (किंवा
ज्यात $\Th{Q}$ चे औपचारिक अर्थनिर्धारण करता येते अशा उपपत्तीत)
संगणनक्षम फलनांचे निरूपण करता येते, हे सहज दिसते. उलट,
ज्या सुसंगत निष्पत्ती व्यवस्थेत भिन्न संख्यांकांची असमता
सिद्ध होते आणि निष्पत्ती ओळखण्याची प्रक्रिया संगणनक्षम
असते, त्या व्यवस्थेत निरूपणीय असलेले प्रत्येक फलन संगणनक्षम
असते. विसंगत व्यवस्थेत प्रत्येक सूत्र सिद्ध होत असल्याने
ही उलटी दिशा तेथे लागू पडत नाही. या दिशेसाठी प्रभावी पुरावा-प्रगणन
पुरेसे आहे. म्हणून, उदाहरणार्थ, या उपपत्ती सुसंगत असतील, तर संगणनक्षम
फलनांचा वर्ग म्हणजे पेआनो अंकगणितात किंवा अगदी
झर्मेलो--फ्रेंकेल संच उपपत्तीत निरूपणीय फलनांचा वर्ग,
असेही लक्षणन करता येते. गोडेल यांनी नोंदवल्याप्रमाणे,
हे काहीसे आश्चर्यकारक आहे. सिद्धतायोग्यतेबाबत प्रश्नांचे
उत्तर कोणती उपपत्ती निवडली आहे यावर फार अवलंबून असते,
हे आपण पुढे पाहू; साधारणपणे स्वयंसिद्धके जितकी बलवान,
तितके अधिक सिद्ध करता येते. पण स्वयंसिद्धकीय उपपत्तींच्या
मोठ्या वर्गात निरूपणीय फलने नेमकी संगणनक्षम फलनेच असतात;
त्या उपपत्ती सुसंगत असतील, भिन्न संख्यांकांची असमता सिद्ध करत असतील
आणि त्यांतील पुराव्यांचे प्रभावी प्रगणन करता येत असेल,
तर अधिक बलवान उपपत्ती अधिक फलनांचे निरूपण करत नाहीत.
\end{explain}












\section{संबंधांचे निरूपण}\label{inc:req:rel:sec}

एखादा \emph{संबंध} निरूपणीय असणे म्हणजे काय, ते आता सांगू.

\begin{defn}
\label{inc:req:rel:defn:representing-relations} नैसर्गिक संख्यांवरील
$R(x_0,\dots,x_k)$ हा संबंध {\em $\Th{Q}$ मध्ये निरूपणीय}
असतो, जर असे सूत्र $\varphi_R(x_0,\dots,x_k)$ असेल की
$R(n_0,\dots,n_k)$ सत्य असेल तेव्हा $\Th{Q}$ हे
$\varphi_R(\num{n_0},\dots,\num{n_k})$ सिद्ध करते, आणि
$R(n_0,\dots,n_k)$ असत्य असेल तेव्हा $\Th{Q}$ हे
$\lnot \varphi_R(\num{n_0},
\dots, \num{n_k})$ सिद्ध करते.
\end{defn}

\begin{thm}
\label{inc:req:rel:thm:representing-rels} संबंध~$\Th{Q}$ मध्ये निरूपणीय
असतो तर आणि तरच तो संगणनक्षम असतो.
\end{thm}

\begin{proof}
प्रथम डावीकडून उजवीकडे जाणारी दिशा पाहू. समजा
$R(x_0,\dots,x_k)$ चे निरूपण $\varphi_R(x_0,\dots,x_k)$ या
सूत्राने होते. $R$ चे संगणन करण्यासाठी पुढील कार्यपद्धती
आहे: $n_0$, \dots,~$n_k$ ही आदाने मिळाल्यावर
$\varphi_R(\num{n_0}, \dots, \num{n_k})$ चा पुरावा आणि
$\lnot \varphi_R(\num{n_0}, \dots, \num{n_k})$ चा पुरावा
यांचा एकाच वेळी शोध घ्या. गृहीतकानुसार यांपैकी एक पुरावा
नक्की सापडेल. पहिला सापडल्यास ``होय'' आणि दुसरा सापडल्यास
``नाही'' असे उत्तर द्या. ही उपपत्ती सुसंगत असल्याने दोन्ही
विरुद्ध पुरावे मिळू शकत नाहीत.

उलट दिशेसाठी $R(x_0, \dots, x_k)$ संगणनक्षम आहे असे समजा.
व्याख्येनुसार याचा अर्थ $\Char{R}(x_0, \dots, x_k)$ हे
फलन संगणनक्षम आहे. \ref{inc:req:int:thm:representable-iff-comp} नुसार
$\Char{R}$ चे एखाद्या सूत्राने, समजा
$\varphi_{\Char{R}}(x_0, \dots, x_k,
y)$ ने, निरूपण होते. $\varphi_R(x_0, \dots, x_k)$ हे
$\varphi_{\Char{R}}(x_0,
\dots, x_k, \num{1})$ सूत्र असू द्या. मग कोणत्याही
$n_0$, \dots,~$n_k$ साठी $R(n_0,
\dots, n_k)$ सत्य असेल, तर $\Char{R}(n_0, \dots, n_k) = 1$.
त्या वेळी $\Th{Q}$ हे
$\varphi_{\Char{R}}(\num{n_0}, \dots, \num{n_k},
\num{1})$ सिद्ध करते; म्हणून $\Th{Q}$ हे
$\varphi_R(\num{n_0}, \dots,
\num{n_k})$ सिद्ध करते. दुसरीकडे $R(n_0, \dots, n_k)$
असत्य असेल, तर $\Char{R}(n_0, \dots, n_k) = 0$.
म्हणून $\Th{Q}$ पुढील विधान सिद्ध करते:
\[\lforall[y][(\varphi_{\Char{R}}(\num{n_0}, \dots, \num{n_k}, y) \lif y =
  \num{0})].
\]
$\Th{Q}$ हे $\eq/[\num{0}][\num{1}]$ सिद्ध करत असल्याने,
$\Th{Q}$ हे $\lnot \varphi_{\Char{R}}(\num{n_0}, \dots, \num{n_k}, \num{1})$
सिद्ध करते; म्हणून $\lnot \varphi_R(\num{n_0}, \dots, \num{n_k})$
हेही सिद्ध करते.
\end{proof}

\begin{prob}
$R$ हा संबंध~$\Th{Q}$ मध्ये निरूपणीय असेल, तर
$\Char{R}$ हेसुद्धा निरूपणीय असते, हे दाखवा.
\end{prob}











\section{अनिर्णेयता}\label{inc:req:und:sec}

एखाद्या उपपत्तीला $\Th{T}$ \emph{अनिर्णेय} म्हणतो, जर $\Th{T}$ च्या
भाषेतील कोणतेही वाक्य~$\varphi$ दिले असता ``$\Th{T}$ हे~$\varphi$
सिद्ध करते का?'' या प्रश्नाचे योग्य उत्तर सांत पायऱ्यांत आणि
नेहमी देणारी संगणकीय पद्धत अस्तित्वात नसेल. म्हणून अंकगणिताच्या
भाषेतील वाक्य~$\varphi$ दिल्यावर $\Th{Q} \Proves \varphi$ आहे की नाही,
हे ठरवणारी संगणकीय पद्धत असेल तर आणि तरच $\Th{Q}$
निर्णेय असेल. हे अधिक नेमके करण्यासाठी विचारू: $\Prov[\Th{Q}](y)$
हा संबंध पुनरावर्ती आहे का? येथे तो~$y$ साठी सत्य
असतो, तर आणि तरच $y$ ही~$\Th{Q}$ मध्ये सिद्ध होणाऱ्या वाक्याची
गोडेल संख्या असते. उत्तर नाही.

\begin{thm}
$\Th{Q}$ अनिर्णेय आहे; म्हणजे पुढील संबंध
\[\Prov[\Th{Q}](y) \defiff \fn{Sent}(y) \land
\lexists[x][\Prf[\Th{Q}](x, y)]
\]
पुनरावर्ती नाही.
\end{thm}

\begin{proof}
तो पुनरावर्ती आहे असे समजा. मग थांबण्याची समस्या अशी
सोडवता येईल: $e$ आणि $n$ दिले असता,
$\cfind{e}(n) \fdefined$ तर आणि तरच असा एखादा~$s$
असतो की $T(e, n, s)$, हे आपल्याला माहीत आहे. येथे $T$ हे
\ref{cmp:rec:nft:thm:kleene-nf} मधील क्लिनीचे विधेय
आहे. $T$ आदिम पुनरावर्ती असल्याने त्याचे~$\Th{Q}$ मध्ये
$\psi_T$ या सूत्राने निरूपण होते; म्हणजे
$\Th{Q}
\Proves \psi_T(\num{e}, \num{n}, \num{s})$ तर आणि तरच
$T(e, n, s)$. जर $\Th{Q}
\Proves \psi_T(\num{e}, \num{n}, \num{s})$ असेल, तर
$ \Th{Q} \Proves
\lexists[y][\psi_T(\num{e}, \num{n}, y)]$ हेसुद्धा मिळते.
असा एकही $s$ नसेल, तर प्रत्येक~$s$ साठी
$\Th{Q} \Proves \lnot \psi_T(\num{e}, \num{n}, \num{s})$.
पण $\Th{Q}$ ही $\omega$-सुसंगत आहे: प्रत्येक~$n \in \Nat$
साठी $\Th{Q}
\Proves \lnot \varphi(\num{n})$ असेल, तर
$\Th{Q}
\Proves/ \lexists[y][\varphi(y)]$. कारण $\Th{Q}$ ची स्वयंसिद्धके
मानक प्रतिमानात~$\Struct{N}$ खरी आहेत. म्हणून
$\Th{Q}
\Proves/ \lexists[y][\psi_T(\num{e}, \num{n}, y)]$.
दुसऱ्या शब्दांत, $\Th{Q}
\Proves \lexists[y][\psi_T(\num{e}, \num{n}, y)]$ तर आणि तरच असा एखादा $s$ असतो की $T(e, n, s)$; म्हणजे तर आणि तरच $\cfind{e}(n) \fdefined$. $e$ आणि~$n$ पासून
$\Gn{\lexists[y][\psi_T(\num{e},
    \num{n}, y)]}$ आपण संगणित करू शकतो; तसे करणारे
$g(e, n)$ हे आदिम पुनरावर्ती फलन असू द्या. मग
\[h(e, n) =
\begin{cases}
1 & \text{जर $\Prov[\Th{Q}](g(e, n))$ असेल}\\
0 & \text{अन्यथा}.
\end{cases}
\]
$\Prov[\Th{Q}]$ पुनरावर्ती असेल, तर यावरून $h$ सुद्धा
पुनरावर्ती ठरेल. पण~$h$ हे
\ref{cmp:rec:hlt:thm:halting-problem} नुसार पुनरावर्ती नाही.
म्हणून $\Prov[\Th{Q}]$ हाही पुनरावर्ती नाही.
\end{proof}

\begin{cor}
प्रथम-क्रम तर्कशास्त्र अनिर्णेय आहे.
\end{cor}

\begin{proof}
प्रथम-क्रम तर्कशास्त्र निर्णेय असते, तर~$\Th{Q}$ मधील
सिद्धतायोग्यताही निर्णेय असती: कारण $\Th{Q} \Proves \varphi$
तर आणि तरच $\Proves \mathsf{T} \lif \varphi$; येथे $\mathsf{T}$ हा
$\Th{Q}$ च्या सांत स्वयंसिद्धकांचा संयोग आहे.
\end{proof}











\section{\texorpdfstring{$\Sigma_1$}{सिग्मा-१} संपूर्णता}\label{inc:inp:s1c:sec}

$\Th{Q}$ आणि तिच्या सुसंगत, स्वयंसिद्धकीकरण करता येणाऱ्या
विस्तारांची अपूर्णता असूनही, संख्यांकांविषयी अनेक मूलभूत तथ्ये
$\Th{Q}$ सिद्ध करते, हे आपण पाहिले आहे. ही क्षमता प्रत्यक्षात
आणखी बरीच व्यापक आहे. $\Th{Q}$ मध्ये काय सिद्ध करता येते
याची व्याप्ती समजण्यासाठी $\Delta_0$, $\Sigma_1$ आणि
$\Pi_1$ सूत्रे यांच्या संकल्पना मांडू. स्थूलमानाने,
$\Sigma_1$ सूत्र हे $\lexists[x][\psi(x)]$ या रूपाचे
असते; येथे $\psi$ फक्त विधानीय संयोजके आणि परिबद्ध संख्यापक
वापरून तयार केलेले असते. $\varphi$ हे $\Struct{N}$ मध्ये सत्य
असलेले $\Sigma_1$ वाक्य असेल, तर
$\Th{Q} \Proves \varphi$, हे आपण दाखवू
(\ref{inc:inp:s1c:thm:sigma1-completeness}).

\begin{defn}
\label{inc:inp:s1c:defn:bd-quant}
\emph{परिबद्ध अस्तित्ववाची सूत्र} हे
$\lexists[x][(x < t \land \varphi(x))]$ या रूपाचे असते; येथे $t$
हे असे कोणतेही पद आहे की त्यात संबंधित संख्यापकाचे चल येत नाही. परंपरेने ते
$\bexists{x < t}{\varphi(x)}$ असे लिहितो.

\emph{परिबद्ध वैश्विक सूत्र} हे
$\lforall[x][(x < t \lif \varphi(x))]$ या रूपाचे असते; येथेही
$t$ या पदात संबंधित संख्यापकाचे चल नसते. परंपरेने ते
$\bforall{x < t}{\varphi(x)}$ असे लिहितो.
\end{defn}

\begin{defn}
\label{inc:inp:s1c:defn:delta0-sigma1-pi1-frm}
सूत्र $\psi$ हे $\Delta_0$ असते, जर ते आण्विक
सूत्रांपासून केवळ विधानीय संयोजके आणि परिबद्ध
संख्यकीकरण वापरून तयार झाले असेल.

सूत्र $\varphi$ हे $\Sigma_1$ असते, जर
$\varphi \ident \lexists[x][\psi(x)]$ आणि $\psi$ हे $\Delta_0$
असेल.

सूत्र $\varphi$ हे $\Pi_1$ असते, जर
$\varphi \ident \lforall[x][\psi(x)]$ आणि $\psi$ हे $\Delta_0$
असेल.
\end{defn}

\begin{lem}
\label{inc:inp:s1c:lem:q-proves-clterm-id} समजा $t$ हे बंद पद असून
$\Value{t}{N} = n$. तर $\Th{Q} \Proves \eq[t][\num n]$.
\end{lem}

\begin{proof}
$t$ च्या गुंतागुंतीवर विगमन करू. आधार प्रकरणात
$\Value{\Obj 0}{N} = 0$, आणि
$\num 0 \ident \Obj 0$ असल्यामुळे
$\Th{Q} \Proves \eq[\Obj 0][\num 0]$.

विगमनाच्या टप्प्यासाठी $t_1$ आणि $t_2$ ही पदे अशी घ्या की
$\Value{t_1}{N} = n_1$, $\Value{t_2}{N} = n_2$,
$\Th{Q} \Proves \eq[t_1][\num n_1]$, आणि
$\Th{Q} \Proves \eq[t_2][\num n_2]$.

मग $\Value{(t_1')}{N} = n_1 + 1$. विगमन गृहीतक आणि
$\eq[\num{n_1}'][\num{n_1}']$ या सूत्र वर
प्रथम-क्रम एकरूपतेचे नियम लावल्यास
$\Th{Q} \Proves \eq[t_1'][{\num n_1}']$ मिळते.
संख्यांकांच्या व्याख्येने
$\Th{Q} \Proves \eq[t_1'][\num{n_1 + 1}]$ मिळते.

बेरीजेसाठी,
\[      \Value{(t_1 + t_2)}{N}
    = \Value{t_1}{N} + \Value{t_2}{N}
    = n_1 + n_2.
\]
विगमन गृहीतक आणि एकरूपतेच्या नियमांनी प्रथम
$\Th{Q} \Proves \eq[t_1 + t_2][\num{n_1} + t_2]$,
आणि ते नियम पुन्हा लावल्यास
$\Th{Q} \Proves \eq[t_1 + t_2][\num{n_1} + \num{n_2}]$
मिळते. \ref{inc:req:bre:lem:q-proves-add} नुसार
$\Th{Q} \Proves \eq[\num{n_1} + \num{n_2}][\num{n_1 + n_2}]$;
म्हणून $\Th{Q} \Proves \eq[t_1 + t_2][\num{n_1 + n_2}]$.

$\times$ साठीही
\ref{inc:req:bre:lem:q-proves-mult} वापरून असाच
युक्तिवाद करता येतो.

अंकगणितातील बंद पदांचे हे सर्व प्रकार आहेत. म्हणून
$\Value{t}{N} = n$ असलेल्या प्रत्येक बंद पद~$t$ साठी
$\Th{Q} \Proves \eq[t][\num n]$.
\end{proof}

\begin{prob}
\ref{inc:inp:s1c:lem:q-proves-clterm-id} मधील $\times$ संबंधीचा
भाग सविस्तर सिद्ध करा.
\end{prob}

\begin{lem}
\label{inc:inp:s1c:lem:atomic-completeness}
समजा $t_1$ आणि $t_2$ ही बंद पदे आहेत. तर
\begin{enumerate}
\item $\Value{t_1}{N} = \Value{t_2}{N}$ असेल, तर
    $\Th{Q} \Proves \eq[t_1][t_2]$.
\item $\Value{t_1}{N} \neq \Value{t_2}{N}$ असेल, तर
    $\Th{Q} \Proves \eq/[t_1][t_2]$.
\item $\Value{t_1}{N} < \Value{t_2}{N}$ असेल, तर
    $\Th{Q} \Proves t_1 < t_2$.
\item $\Value{t_2}{N} \leq \Value{t_1}{N}$ असेल, तर
    $\Th{Q} \Proves \lnot(t_1 < t_2)$.
\end{enumerate}
\end{lem}

\begin{proof}
दिलेले $t_1$ आणि $t_2$ घेऊन
$n = \Value{t_1}{N}$ आणि $m = \Value{t_2}{N}$
निश्चित करू.

समजा $\varphi \ident t_1 = t_2$.
\ref{inc:inp:s1c:lem:q-proves-clterm-id} नुसार
$\Th{Q} \Proves \eq[t_1][\num n]$ आणि
$\Th{Q} \Proves \eq[t_2][\num m]$.
$n = m$ असेल, तर
$\Th{Q} \Proves \eq[\num n][\num m]$;
म्हणून एकरूपतेच्या संक्रमणशीलतेने
$\Th{Q} \Proves \eq[t_1][t_2]$.
$n \neq m$ असेल, तर
$\Th{Q} \Proves \eq/[\num n][\num m]$;
एकरूपतेचे नियम पुन्हा वापरल्यास
$\Th{Q} \Proves \eq/[t_1][t_2]$.

आता $\varphi \ident t_1 < t_2$ असे घेऊ. दोन्ही प्रकरणांत
$\mathsf{Q}_8$ या स्वयंसिद्धकाचा आधार घेऊ. त्यानुसार सर्व
$x,y$ साठी
$x < y \liff \lexists[z][\eq[z' + x][y]]$.

समजा $\Sat{N}{t_1 < t_2}$. मग असा $k \in \Nat$
असतो की $n + k + 1 = m$.
\ref{inc:inp:s1c:lem:q-proves-clterm-id} नुसार
$\Th{Q} \Proves \eq[t_1][\num n]$ आणि
$\Th{Q} \Proves \eq[t_2][\num m]$.
या प्रमेयाच्या पहिल्या भागाने निश्चित संख्यांकांची बेरीज
काढल्यास
$\Th{Q} \Proves \eq[{\num k}' + \num n][\num m]$.
एकरूपतेच्या संक्रमणशीलतेने मग
$\Th{Q} \Proves \eq[{\num k}' + t_1][t_2]$;
म्हणून
$\Th{Q} \Proves \lexists[z][\eq[z' + t_1][t_2]]$.
$\mathsf{Q}_8$ ची उजवीकडून डावीकडे दिशा वापरल्यास
$\Th{Q} \Proves t_1 < t_2$.

याउलट $\Sat/{N}{t_1 < t_2}$, म्हणजे $m \leq n$,
असे समजा.

$\Th{Q}$ मध्ये काम करून $t_1 < t_2$ असे गृहीत धरा.
$\mathsf{Q}_8$ ची डावीकडून उजवीकडे दिशा लावल्यास असा $z$
मिळतो की $\eq[z' + t_1][t_2]$.
$\Th{Q} \Proves \eq[t_1][\num n]$ आणि
$\Th{Q} \Proves \eq[t_2][\num m]$ असल्यामुळे
$\eq[z' + \num n][\num m]$.

$m$ वर बाह्य विगमन करू: प्रत्येक पायरीवर $\mathsf{Q}_5$
वापरून बेरीज उत्तरवर्तीच्या रूपात लिहू आणि उत्तरवर्तीचे
एकास-एक असणे वापरून दोन्ही बाजूंचे एक-एक उत्तरवर्ती
रद्द करू. अशा रीतीने
$\eq[z' + \num{n - m}][\Obj 0]$ मिळते.
$m = n$ असेल, तर शून्य जोडण्याच्या स्वयंसिद्धकाने
$\eq[z'][\Obj 0]$ मिळते; पण $\mathsf{Q}_2$ नुसार उत्तरवर्ती
शून्य नसतो, म्हणून विरोधाभास मिळतो.
$m < n$ असेल, तर $\mathsf{Q}_5$ पुन्हा वापरून
$\eq[(z' + \num{n - m - 1})'][\Obj 0]$ मिळते;
$\mathsf{Q}_2$ मुळे पुन्हा विरोधाभास मिळतो.
म्हणून $\Th{Q} \Proves \lnot(t_1 < t_2)$.
\end{proof}

\begin{lem}
\label{inc:inp:s1c:lem:bounded-quant-equiv}
समजा $\varphi$ हे सूत्र, $t$ हे बंद पद, आणि
$k=\Value{t}{N}$. तर
\begin{enumerate}
\item $\Th{Q} \Proves \bforall{x<t}{\varphi(x)}$ तर आणि तरच $\Th{Q} \Proves
    \varphi(\num 0) \land \dots \land \varphi(\num{k-1})$.
\item $\Th{Q} \Proves \bexists{x<t}{\varphi(x)}$ तर आणि तरच $\Th{Q} \Proves
    \varphi(\num 0) \lor \dots \lor \varphi(\num{k-1})$.
\end{enumerate}
\end{lem}

\begin{proof}
    परिबद्ध वैश्विक संख्यापकाचे प्रकरण सिद्ध करू.
    $\Value{t}{N} = 0$ असेल, तर
    \ref{inc:req:min:lem:less-zero} नुसार
    $x<\num 0$ अशी कोणतीही संख्या नसते; म्हणून
    समतुल्यतेची डावी बाजू $\Th{Q}$ मध्ये सिद्ध होते.
    उजवीकडील रिक्त संयोगाला $\ltrue$ मानतो;
    तेही $\Th{Q}$ मध्ये सिद्ध होते.

    म्हणून आता एखाद्या नैसर्गिक संख्या~$k$ साठी
    $\Value{t}{N} = k+1$ असे समजू.
    \ref{inc:inp:s1c:lem:q-proves-clterm-id} नुसार
    $\bforall{x<\num{k+1}}{\varphi(x)}$ या रूपाच्या
    सूत्र साठी काम करणे पुरेसे आहे.

    समजा
    $\Th{Q} \Proves \bforall{x<\num{k+1}}{\varphi(x)}$,
    आणि $n \leq k$ घेऊ.
    \ref{inc:inp:s1c:lem:atomic-completeness} नुसार
    $\Th{Q} \Proves \num n < \num{k+1}$;
    म्हणून तर्कनियमांनी
    $\Th{Q} \Proves \varphi(\num n)$.
    प्रत्येक $n \leq k$ साठी हे $k+1$ वेळा
    लावल्यास अपेक्षेप्रमाणे
    $\Th{Q} \Proves \varphi(\num 0)
    \land \dots \land \varphi(\num k)$ मिळते.

    उलट दिशेसाठी समजा
    $\Th{Q} \Proves
    \varphi(\num 0) \land \dots \land \varphi(\num k)$.
    $\Th{Q}$ मध्ये $x < \num{k+1}$ असे समजू.
    \ref{inc:req:min:lem:less-nsucc} नुसार
    $x = \num 0 \lor \dots \lor x = \num k$.
    मग तर्कनियमांनी $\varphi(x)$ मिळते आणि त्यावरून
    $\bforall{x<\num{k+1}}{\varphi(x)}$ हे वैश्विक विधान
    मिळते.

    परिबद्ध अस्तित्ववाची संख्यकीकरण असलेल्या
    सूत्रे साठी समतुल्यतेचा पुरावाही असाच आहे.
\end{proof}

\begin{prob}
\ref{inc:inp:s1c:lem:bounded-quant-equiv} मधील अस्तित्ववाची
प्रकरणाचा सविस्तर पुरावा द्या.
\end{prob}

\begin{lem}
\label{inc:inp:s1c:lem:delta0-completeness}
$\varphi$ हे $\Struct{N}$ मध्ये सत्य असलेले
$\Delta_0$ वाक्य असेल, तर
$\Th{Q} \Proves \varphi$.
\end{lem}

\begin{proof}
सूत्राच्या गुंतागुंतीवर विगमन करू.

आधार प्रकरण \ref{inc:inp:s1c:lem:atomic-completeness} ने
मिळते; आता विगमनाचा टप्पा पाहू. सोयीसाठी,
निषेधाच्या प्रकरणाची विभागणी निषिद्ध सूत्र
च्या रचनेनुसार करू. इतर विधानीय संयोजक पुढील
प्रकरणांमधून संक्षेपाने व्यक्त करता येतात.

\begin{enumerate}
\item $(\varphi \land \psi)$ हे $\Struct{N}$ मध्ये सत्य
असेल, तर $\varphi$ आणि $\psi$ दोन्ही $\Struct{N}$ मध्ये सत्य आहेत.
विगमन गृहीतकाने
$\Th{Q} \Proves \varphi$ आणि
$\Th{Q} \Proves \psi$; म्हणून तर्कनियमांनी
$\Th{Q} \Proves (\varphi \land \psi)$.

\item $\lnot (\varphi \land \psi)$ हे $\Struct{N}$
मध्ये सत्य असेल, तर $\lnot \varphi$ किंवा
$\lnot \psi$ हे $\Struct{N}$ मध्ये सत्य आहे. सामान्यत्वाला बाधा
न आणता पहिले सत्य आहे असे धरू.
विगमन गृहीतकाने $\Th{Q} \Proves \lnot \varphi$;
म्हणून तर्कनियमांनी
$\Th{Q} \Proves \lnot (\varphi \land \psi)$.

\item $(\varphi \lor \psi)$ हे $\Struct{N}$ मध्ये सत्य
असेल, तर $\varphi$ हे $\Struct{N}$ मध्ये सत्य आहे किंवा
$\psi$ हे $\Struct{N}$ मध्ये सत्य आहे.
सामान्यत्वाला बाधा न आणता पहिले सत्य आहे
असे धरू. विगमन गृहीतकाने
$\Th{Q} \Proves \varphi$;
म्हणून तर्कनियमांनी
$\Th{Q} \Proves (\varphi \lor \psi)$.

\item $\lnot(\varphi \lor \psi)$ हे $\Struct{N}$
मध्ये सत्य असेल, तर $\lnot \varphi$ आणि
$\lnot \psi$ दोन्ही $\Struct{N}$ मध्ये सत्य आहेत.
विगमन गृहीतकाने
$\Th{Q} \Proves \lnot \varphi$ आणि
$\Th{Q} \Proves \lnot \psi$.
त्यामुळे तर्कनियमांनी
$\Th{Q} \Proves \lnot(\varphi \lor \psi)$.

\item $\bforall{x<t}{\varphi(x)}$ हे
$\Struct{N}$ मध्ये सत्य आहे असे समजा;
येथे $t$ हे बंद पद आणि $k=\Value{t}{N}$.
सर्व $n < \Value{t}{N}$ साठी
$\varphi(\num n)$ हे $\Struct{N}$ मध्ये सत्य
असेल, तर विगमन गृहीतक आणि तर्कनियमांनी
$\Th{Q} \Proves
\varphi(\num 0) \land \dots \land \varphi(\num{k-1})$.
\ref{inc:inp:s1c:lem:bounded-quant-equiv} नुसार
$\Th{Q} \Proves \bforall{x<t}{\varphi(x)}$.

\item परिबद्ध अस्तित्ववाची संख्यापकाचे
प्रकरणही वैश्विक संख्यापकासारखेच आहे:
येथे $\bexists{x < t}{\varphi(x)}$ या रूपाचे
वाक्य असते.

\item $\lnot \bforall{x<t}{\varphi(x)}$ हे
$\Struct{N}$ मध्ये सत्य आहे असे समजा;
येथे $t$ हे बंद पद. हे वाक्य
$\bexists{x<t}{\lnot \varphi(x)}$ या वाक्य
शी समतुल्य आहे आणि ही समतुल्यता
$\Th{Q}$ मध्येही सिद्ध करता येते. म्हणून
परिबद्ध अस्तित्ववाची संख्यापकाचा युक्तिवाद
लागू करता येतो.

\item त्याचप्रमाणे, $t$ हे बंद पद असताना
$\lnot \bexists{x<t}{\varphi(x)}$ हे
$\Struct{N}$ मध्ये सत्य आहे असे समजा.
हे वाक्य $\Th{Q}$ मध्ये
$\bforall{x<t}{\lnot\varphi(x)}$ शी समतुल्य
आहे; म्हणून परिबद्ध वैश्विक संख्यापकाचा
युक्तिवाद लागू करता येतो.

\item अखेरीस $\lnot \varphi$ हे
$\Struct{N}$ मध्ये सत्य आहे असे समजा.
उरलेली प्रकरणे अशी: $\varphi$ आण्विक आहे,
किंवा एखाद्या $\Delta_0$ वाक्य $\psi$
साठी $\lnot \varphi \ident \lnot\lnot \psi$.
$\varphi$ आण्विक असेल, तर
\ref{inc:inp:s1c:lem:atomic-completeness} नुसार
$\Th{Q} \Proves \lnot \varphi$.
$\lnot \varphi \ident \lnot\lnot \psi$ असेल,
तर तर्कनियमांनी ते $\Th{Q}$ मध्ये
$\psi$ शी सिद्धतया समतुल्य आहे.
$\lnot \varphi$ हे $\Struct{N}$ मध्ये सत्य
असल्याने तेही $\Struct{N}$ मध्ये सत्य आहे.
म्हणून विगमन गृहीतकाने
$\Th{Q} \Proves \lnot \varphi$.
\end{enumerate}
\end{proof}

\begin{prob}
\ref{inc:inp:s1c:lem:delta0-completeness} मधील
अस्तित्ववाची प्रकरणाचा सविस्तर पुरावा द्या.
\end{prob}

\begin{thm}
\label{inc:inp:s1c:thm:sigma1-completeness}
$\varphi$ हे $\Struct{N}$ मध्ये सत्य असलेले
$\Sigma_1$ वाक्य असेल, तर
$\Th{Q} \Proves \varphi$.
\end{thm}

\begin{proof}
$\lexists[x][\varphi(x)]$ हे $\Struct{N}$ मध्ये
सत्य असलेले $\Sigma_1$ वाक्य असेल,
तर अशी नैसर्गिक संख्या~$n$ आणि चल-मूल्यन~$s$
असतात की $s(x) = n$ आणि
$\Sat{N}{\varphi(x)}[s]$.
पूर्ति-संबंधाच्या नेहमीच्या गुणधर्मांनुसार
$\Sat{N}{\varphi(\num n)}$.
परंतु $\varphi(\num n)$ हे बंद $\Delta_0$ सूत्र
आहे; म्हणून \ref{inc:inp:s1c:lem:delta0-completeness} ने
$\Th{Q} \Proves \varphi(\num n)$.
मग तर्कनियमांनी
$\Th{Q} \Proves \lexists[x][\varphi(x)]$.
\end{proof}



